% Options for packages loaded elsewhere
\PassOptionsToPackage{unicode}{hyperref}
\PassOptionsToPackage{hyphens}{url}
\documentclass[
  10pt,
  a4paper,
  oneside]{report}
\usepackage{xcolor}
\usepackage[margin=18mm]{geometry}
\usepackage{amsmath,amssymb}
\setcounter{secnumdepth}{5}
\usepackage{iftex}
\ifPDFTeX
  \usepackage[T1]{fontenc}
  \usepackage[utf8]{inputenc}
  \usepackage{textcomp} % provide euro and other symbols
\else % if luatex or xetex
  \usepackage{unicode-math} % this also loads fontspec
  \defaultfontfeatures{Scale=MatchLowercase}
  \defaultfontfeatures[\rmfamily]{Ligatures=TeX,Scale=1}
\fi
\usepackage{lmodern}
\ifPDFTeX\else
  % xetex/luatex font selection
\fi
% Use upquote if available, for straight quotes in verbatim environments
\IfFileExists{upquote.sty}{\usepackage{upquote}}{}
\IfFileExists{microtype.sty}{% use microtype if available
  \usepackage[]{microtype}
  \UseMicrotypeSet[protrusion]{basicmath} % disable protrusion for tt fonts
}{}
\usepackage{setspace}
\makeatletter
\@ifundefined{KOMAClassName}{% if non-KOMA class
  \IfFileExists{parskip.sty}{%
    \usepackage{parskip}
  }{% else
    \setlength{\parindent}{0pt}
    \setlength{\parskip}{6pt plus 2pt minus 1pt}}
}{% if KOMA class
  \KOMAoptions{parskip=half}}
\makeatother
\usepackage{color}
\usepackage{fancyvrb}
\newcommand{\VerbBar}{|}
\newcommand{\VERB}{\Verb[commandchars=\\\{\}]}
\DefineVerbatimEnvironment{Highlighting}{Verbatim}{commandchars=\\\{\}}
% Add ',fontsize=\small' for more characters per line
\newenvironment{Shaded}{}{}
\newcommand{\AlertTok}[1]{\textcolor[rgb]{1.00,0.00,0.00}{\textbf{#1}}}
\newcommand{\AnnotationTok}[1]{\textcolor[rgb]{0.38,0.63,0.69}{\textbf{\textit{#1}}}}
\newcommand{\AttributeTok}[1]{\textcolor[rgb]{0.49,0.56,0.16}{#1}}
\newcommand{\BaseNTok}[1]{\textcolor[rgb]{0.25,0.63,0.44}{#1}}
\newcommand{\BuiltInTok}[1]{\textcolor[rgb]{0.00,0.50,0.00}{#1}}
\newcommand{\CharTok}[1]{\textcolor[rgb]{0.25,0.44,0.63}{#1}}
\newcommand{\CommentTok}[1]{\textcolor[rgb]{0.38,0.63,0.69}{\textit{#1}}}
\newcommand{\CommentVarTok}[1]{\textcolor[rgb]{0.38,0.63,0.69}{\textbf{\textit{#1}}}}
\newcommand{\ConstantTok}[1]{\textcolor[rgb]{0.53,0.00,0.00}{#1}}
\newcommand{\ControlFlowTok}[1]{\textcolor[rgb]{0.00,0.44,0.13}{\textbf{#1}}}
\newcommand{\DataTypeTok}[1]{\textcolor[rgb]{0.56,0.13,0.00}{#1}}
\newcommand{\DecValTok}[1]{\textcolor[rgb]{0.25,0.63,0.44}{#1}}
\newcommand{\DocumentationTok}[1]{\textcolor[rgb]{0.73,0.13,0.13}{\textit{#1}}}
\newcommand{\ErrorTok}[1]{\textcolor[rgb]{1.00,0.00,0.00}{\textbf{#1}}}
\newcommand{\ExtensionTok}[1]{#1}
\newcommand{\FloatTok}[1]{\textcolor[rgb]{0.25,0.63,0.44}{#1}}
\newcommand{\FunctionTok}[1]{\textcolor[rgb]{0.02,0.16,0.49}{#1}}
\newcommand{\ImportTok}[1]{\textcolor[rgb]{0.00,0.50,0.00}{\textbf{#1}}}
\newcommand{\InformationTok}[1]{\textcolor[rgb]{0.38,0.63,0.69}{\textbf{\textit{#1}}}}
\newcommand{\KeywordTok}[1]{\textcolor[rgb]{0.00,0.44,0.13}{\textbf{#1}}}
\newcommand{\NormalTok}[1]{#1}
\newcommand{\OperatorTok}[1]{\textcolor[rgb]{0.40,0.40,0.40}{#1}}
\newcommand{\OtherTok}[1]{\textcolor[rgb]{0.00,0.44,0.13}{#1}}
\newcommand{\PreprocessorTok}[1]{\textcolor[rgb]{0.74,0.48,0.00}{#1}}
\newcommand{\RegionMarkerTok}[1]{#1}
\newcommand{\SpecialCharTok}[1]{\textcolor[rgb]{0.25,0.44,0.63}{#1}}
\newcommand{\SpecialStringTok}[1]{\textcolor[rgb]{0.73,0.40,0.53}{#1}}
\newcommand{\StringTok}[1]{\textcolor[rgb]{0.25,0.44,0.63}{#1}}
\newcommand{\VariableTok}[1]{\textcolor[rgb]{0.10,0.09,0.49}{#1}}
\newcommand{\VerbatimStringTok}[1]{\textcolor[rgb]{0.25,0.44,0.63}{#1}}
\newcommand{\WarningTok}[1]{\textcolor[rgb]{0.38,0.63,0.69}{\textbf{\textit{#1}}}}
\usepackage{longtable,booktabs,array}
\usepackage{caption}
\captionsetup[table]{skip=6pt}
\newcounter{none} % for unnumbered tables
\usepackage{calc} % for calculating minipage widths
% Correct order of tables after \paragraph or \subparagraph
\usepackage{etoolbox}
\makeatletter
\patchcmd\longtable{\par}{\if@noskipsec\mbox{}\fi\par}{}{}
\makeatother
% Allow footnotes in longtable head/foot
\IfFileExists{footnotehyper.sty}{\usepackage{footnotehyper}}{\usepackage{footnote}}
\makesavenoteenv{longtable}
\setlength{\emergencystretch}{3em} % prevent overfull lines
\providecommand{\tightlist}{%
  \setlength{\itemsep}{0pt}\setlength{\parskip}{0pt}}
\usepackage{fancyhdr}
\usepackage{microtype}
\usepackage{enumitem}
\usepackage{array}
\usepackage{xurl}
\definecolor{UnitelBlue}{HTML}{2457A6}
\definecolor{UnitelDark}{HTML}{17324D}
\setlength{\emergencystretch}{3em}
\setlength{\headheight}{15pt}
\renewcommand{\arraystretch}{1.14}
\setlength{\tabcolsep}{4pt}
\Urlmuskip=0mu plus 1mu
\setlist[itemize]{itemsep=2pt,topsep=4pt}
\setlist[enumerate]{itemsep=2pt,topsep=4pt}
\raggedbottom
\pagestyle{fancy}
\fancyhf{}
\fancyhead[L]{\small\textcolor{UnitelDark}{UNI-TEL Software Engineering Project}}
\fancyhead[R]{\small\textcolor{UnitelDark}{Five-Day Master Plan}}
\fancyfoot[C]{\small\thepage}
\renewcommand{\headrulewidth}{0.4pt}
\renewcommand{\footrulewidth}{0pt}
\PassOptionsToPackage{colorlinks=true,linkcolor=UnitelBlue,urlcolor=UnitelBlue,citecolor=UnitelBlue}{hyperref}
\usepackage{bookmark}
\IfFileExists{xurl.sty}{\usepackage{xurl}}{} % add URL line breaks if available
\urlstyle{same}
% fallback for those not using the hyperref driver hyperxmp:
\makeatletter
\@ifundefined{xmpquote}{\newcommand{\xmpquote}[1]{#1}}{}
\makeatother
\hypersetup{
  pdftitle={UNI-TEL Five-Day Software Engineering Project Master Plan},
  pdfauthor={UNI-TEL Project Team},
  hidelinks,
  pdfcreator={LaTeX via pandoc}}

\title{UNI-TEL Five-Day Software Engineering Project Master Plan}
\usepackage{etoolbox}
\makeatletter
\providecommand{\subtitle}[1]{% add subtitle to \maketitle
  \apptocmd{\@title}{\par {\large #1 \par}}{}{}
}
\makeatother
\subtitle{Five-member execution plan for 14-18 August 2026}
\author{UNI-TEL Project Team}
\date{Version 1.0 - 14 August 2026}

\begin{document}
\maketitle

{
\setcounter{tocdepth}{1}
\tableofcontents
}
\setstretch{1.04}
\begin{quote}
\textbf{For the project team:} Execute this plan through GitHub issues
and pull requests. Track every checkbox, attach genuine evidence, and do
not mark an activity complete until its stated exit criteria are
satisfied.
\end{quote}

\textbf{Goal:} Convert the imported UNI-TEL prototype into a traceable,
tested, documented, and presentable Software Engineering class project
in five calendar days with five team members.

\textbf{Approach:} Use Rapid Application Development (RAD) as the
primary lifecycle, iterative and incremental delivery for daily
integration, Kanban for task flow, prototyping for UI validation,
V-Model-style requirement-to-test traceability, and DevOps practices for
automated quality checks and release management.

\textbf{Technology baseline:} React 18, TypeScript, Vite, Supabase
Authentication and PostgreSQL, Tailwind CSS, shadcn/ui, TanStack Query,
Recharts, Git, GitHub, GitHub Actions, Vitest, React Testing Library,
Vercel, Markdown, LaTeX.

\textbf{Official execution window:} Friday, 14 August 2026 through
Tuesday, 18 August 2026 (IST).

\begin{center}\rule{0.5\linewidth}{0.5pt}\end{center}

\chapter{Document Control}\label{1-document-control}

{\def\LTcaptype{none} % do not increment counter
\begin{longtable}[]{@{}
  >{\raggedright\arraybackslash}p{(\linewidth - 2\tabcolsep) * \real{0.3333}}
  >{\raggedright\arraybackslash}p{(\linewidth - 2\tabcolsep) * \real{0.6667}}@{}}
\toprule\noalign{}
\begin{minipage}[b]{\linewidth}\raggedright
Field
\end{minipage} & \begin{minipage}[b]{\linewidth}\raggedright
Value
\end{minipage} \\
\midrule\noalign{}
\endhead
\bottomrule\noalign{}
\endlastfoot
Project & UNI-TEL - Student Academic Management and Analytics
Platform \\
Document & Five-Day Software Engineering Master Plan \\
Version & 1.0 \\
Planning date & 14 August 2026 \\
Duration & 5 calendar days \\
Team size & 5 members \\
Primary methodology & Rapid Application Development \\
Supporting practices & Iterative delivery, Kanban, prototyping, V-Model
traceability, DevOps \\
Repository baseline & Existing UNI-TEL code imported as one clearly
identified baseline commit \\
Final release target & \texttt{v1.0.0-se-class} \\
\end{longtable}
}

Replace the labels Member 1 through Member 5 with actual names in the
Project Charter on Day 1. Keep the role labels in addition to names so
ownership remains clear.

\chapter{How to Use This Plan}\label{2-how-to-use-this-plan}

\begin{enumerate}
\def\labelenumi{\arabic{enumi}.}
\tightlist
\item
  Create the new GitHub repository and import the existing application
  as one baseline commit.
\item
  Create one GitHub milestone named
  \texttt{UNI-TEL\ SE\ Class\ -\ 18\ Aug\ 2026}.
\item
  Create issues from the work-package and backlog tables in this plan.
\item
  Assign one accountable owner and one reviewer to each issue.
\item
  Move issues across the Kanban board as evidence is produced.
\item
  Use feature, documentation, test, fix, or CI branches. Do not work
  directly on \texttt{main}.
\item
  Submit each coherent change through a pull request linked to its issue
  and requirement IDs.
\item
  Run the daily review gate before ending each day.
\item
  Update the requirements traceability matrix continuously, not only on
  Day 5.
\item
  Complete the final release gate before tagging the submission.
\end{enumerate}

\chapter{Project Integrity and Baseline
Policy}\label{3-project-integrity-and-baseline-policy}

The new repository represents the official five-day class effort. The
application code already exists, so its origin must be handled
transparently.

\section{Required baseline action}\label{31-required-baseline-action}

The first repository commit must be:

\begin{Shaded}
\begin{Highlighting}[]
\NormalTok{chore: import UNI{-}TEL baseline prototype}
\end{Highlighting}
\end{Shaded}

The baseline commit may contain the current application code, but it
must exclude the former \texttt{.git} directory, \texttt{node\_modules},
\texttt{.env}, credentials, generated build output, local editor state,
and temporary files.

\section{Required baseline note}\label{32-required-baseline-note}

Add this statement to the new repository README:

\begin{quote}
This repository began with an imported UNI-TEL application prototype.
During the official five-day Software Engineering project window, the
team formalized requirements and design, established traceability, added
testing and quality controls, corrected verified defects, documented
deployment and maintenance, and prepared the assessed release.
\end{quote}

\section{Evidence rules}\label{33-evidence-rules}

\begin{itemize}
\tightlist
\item
  Do not backdate commits, forge interviews, invent user feedback, or
  manufacture test results.
\item
  Do not split old source files into artificial commits to simulate
  implementation history.
\item
  Record only work genuinely completed during the five-day period.
\item
  Mark a test Pass only after execution and preservation of its result.
\item
  Mark a review complete only after another member provides a
  substantive review.
\item
  Preserve citations and licenses for reused libraries, components,
  templates, and assets.
\item
  Distinguish measured values from targets. For example, a performance
  goal is not a measured result until a test report exists.
\end{itemize}

\chapter{Project Definition}\label{4-project-definition}

\section{Problem statement}\label{41-problem-statement}

Students often manage semesters, subjects, marks, attendance, and
academic progress using disconnected notebooks, spreadsheets, portals,
or calculators. Fragmented data makes it difficult to calculate SGPA and
CGPA consistently, identify attendance shortages, evaluate trends, and
generate useful reports. UNI-TEL provides one secure, responsive
platform in which an individual student can maintain and analyze
academic records.

\section{Product vision}\label{42-product-vision}

UNI-TEL will be a student-centered academic tracking application that
converts raw marks and attendance data into understandable progress
information while preserving user ownership, data integrity, and
portability.

\section{Objectives}\label{43-objectives}

\begin{itemize}
\tightlist
\item
  Centralize semester, subject, marks, and attendance information.
\item
  Calculate attendance percentage, SGPA, CGPA, and weighted marks
  consistently.
\item
  Help students recognize weak attendance and performance trends.
\item
  Provide analytics and exportable academic reports.
\item
  Isolate each user\textquotesingle s data through authentication and
  database security.
\item
  Demonstrate a complete and traceable Software Engineering lifecycle.
\item
  Deliver a reproducible build, tested release, user guide, and
  maintenance plan.
\end{itemize}

\section{Stakeholders}\label{44-stakeholders}

{\def\LTcaptype{none} % do not increment counter
\begin{longtable}[]{@{}
  >{\raggedright\arraybackslash}p{(\linewidth - 6\tabcolsep) * \real{0.2393}}
  >{\raggedright\arraybackslash}p{(\linewidth - 6\tabcolsep) * \real{0.3419}}
  >{\raggedleft\arraybackslash}p{(\linewidth - 6\tabcolsep) * \real{0.0769}}
  >{\raggedright\arraybackslash}p{(\linewidth - 6\tabcolsep) * \real{0.3419}}@{}}
\toprule\noalign{}
\begin{minipage}[b]{\linewidth}\raggedright
Stakeholder
\end{minipage} & \begin{minipage}[b]{\linewidth}\raggedright
Interest
\end{minipage} & \begin{minipage}[b]{\linewidth}\raggedleft
Influence
\end{minipage} & \begin{minipage}[b]{\linewidth}\raggedright
Engagement
\end{minipage} \\
\midrule\noalign{}
\endhead
\bottomrule\noalign{}
\endlastfoot
Student user & Accurate, private, usable academic tracking & High &
Primary requirements and acceptance viewpoint \\
Course professor & Evidence of SE theory and disciplined execution &
High & Reviews final artifacts and demonstration \\
Five-member project team & Deliverable quality and fair contribution &
High & Daily planning, implementation, review, presentation \\
Supabase & Backend platform and service constraints & Medium &
Architecture and operational dependency \\
Vercel or selected host & Frontend deployment & Medium & Deployment and
release dependency \\
Future teachers/institutions & Possible multi-user extension & Low &
Future scope only \\
\end{longtable}
}

\section{In-scope product
capabilities}\label{45-in-scope-product-capabilities}

\begin{itemize}
\tightlist
\item
  Registration, login, logout, session handling, and protected routes.
\item
  Student profile management.
\item
  Semester creation, update, viewing, and deletion.
\item
  Subject and credit management.
\item
  Attendance recording and percentage calculation.
\item
  Examination marks, total marks, weightage, and percentage management.
\item
  SGPA and CGPA calculation.
\item
  Performance dashboards, trends, and academic analytics.
\item
  JSON import/export and PDF/Excel reporting where the baseline supports
  them.
\item
  User notifications where the baseline supports them.
\item
  Responsive web interface, validation, error handling, and
  empty/loading states.
\item
  Supabase schema, constraints, migrations, authentication, and
  Row-Level Security.
\item
  Automated quality checks, test evidence, deployment documentation, and
  release packaging.
\end{itemize}

\section{Explicitly out of scope}\label{46-explicitly-out-of-scope}

\begin{itemize}
\tightlist
\item
  Teacher, administrator, and institution-wide dashboards.
\item
  Knowledge Hub or study-material sharing.
\item
  Social networking, groups, leaderboards, or chat.
\item
  AI recommendations or grade prediction presented as production AI.
\item
  Native Android or iOS applications.
\item
  Payment processing, LMS integration, and bulk institutional imports.
\item
  A complete redesign of the existing application.
\item
  Unverified claims of production scale, guaranteed uptime, WCAG
  certification, or benchmark performance.
\end{itemize}

\section{Constraints and
assumptions}\label{47-constraints-and-assumptions}

\begin{itemize}
\tightlist
\item
  The project has five calendar days and five contributors.
\item
  The imported prototype is the technical baseline.
\item
  Internet access is required for hosted Supabase and deployment
  services.
\item
  The team will use test or synthetic academic data only.
\item
  The final system is an academic prototype, not a certified production
  service.
\item
  Scope changes require a recorded change request and Project Lead
  approval.
\item
  \texttt{supabase/migrations} is the authoritative schema source until
  the team verifies otherwise.
\end{itemize}

\chapter{Success Criteria}\label{5-success-criteria}

The project is successful only when all of the following are true:

\begin{enumerate}
\def\labelenumi{\arabic{enumi}.}
\tightlist
\item
  The scope, objectives, stakeholders, feasibility, and methodology are
  documented.
\item
  Every approved functional and non-functional requirement has a stable
  identifier.
\item
  Major requirements have use cases and measurable acceptance criteria.
\item
  Architecture, behavior, database, component, and deployment views are
  documented consistently.
\item
  Each Must-have requirement maps to implementation and at least one
  executed test in the RTM.
\item
  Critical calculations and validation rules have automated tests.
\item
  Authentication, authorization, data integrity, error handling, and
  responsive behavior have test evidence.
\item
  CI runs installation, type checking, linting, automated tests, and
  production build.
\item
  All Critical and High release-blocking defects are closed or formally
  accepted with justification.
\item
  The application builds reproducibly and the deployment procedure is
  documented.
\item
  PRD, SRS, design, test, deployment, maintenance, user, and
  final-report artifacts agree with one another.
\item
  Every team member has attributable issues, commits, pull requests,
  reviews, technical work, documentation work, and a presentation
  segment.
\item
  The final release is tagged \texttt{v1.0.0-se-class} and accompanied
  by release notes.
\item
  The presentation demonstrates theory through direct project evidence,
  not definitions alone.
\end{enumerate}

\chapter{Methodology Coverage and
Selection}\label{6-methodology-coverage-and-selection}

\section{Methodologies to explain in the
report}\label{61-methodologies-to-explain-in-the-report}

{\def\LTcaptype{none} % do not increment counter
\begin{longtable}[]{@{}
  >{\raggedright\arraybackslash}p{(\linewidth - 8\tabcolsep) * \real{0.1304}}
  >{\raggedright\arraybackslash}p{(\linewidth - 8\tabcolsep) * \real{0.2174}}
  >{\raggedright\arraybackslash}p{(\linewidth - 8\tabcolsep) * \real{0.2174}}
  >{\raggedright\arraybackslash}p{(\linewidth - 8\tabcolsep) * \real{0.2174}}
  >{\raggedright\arraybackslash}p{(\linewidth - 8\tabcolsep) * \real{0.2174}}@{}}
\toprule\noalign{}
\begin{minipage}[b]{\linewidth}\raggedright
Method or model
\end{minipage} & \begin{minipage}[b]{\linewidth}\raggedright
Core idea
\end{minipage} & \begin{minipage}[b]{\linewidth}\raggedright
Strength
\end{minipage} & \begin{minipage}[b]{\linewidth}\raggedright
Limitation
\end{minipage} & \begin{minipage}[b]{\linewidth}\raggedright
UNI-TEL decision
\end{minipage} \\
\midrule\noalign{}
\endhead
\bottomrule\noalign{}
\endlastfoot
Waterfall & Complete sequential phases & Clear documents and phase gates
& Change is expensive and feedback arrives late & Explain for theory; do
not claim as primary \\
V-Model & Pair development stages with test levels & Strong verification
and validation traceability & Rigid when requirements change & Use
requirement-to-test mapping only \\
Prototyping & Learn through early UI or behavior models & Fast feedback
and requirement clarification & Prototype may be mistaken for production
quality & Use the imported UI as the starting prototype \\
Incremental & Deliver functionality in usable slices & Early value and
easier integration & Requires disciplined interfaces & Use daily
integrated increments \\
Iterative & Revisit and improve earlier work & Supports discovery and
correction & Can cause uncontrolled rework & Use daily review and
refinement \\
Spiral & Risk-driven cycles & Strong for large, uncertain, high-risk
systems & Too costly and complex for five days & Explain and reject for
this scope \\
RAD & Time-boxed construction with reuse and prototyping & Fits short
deadlines and existing components & Documentation and architecture can
be neglected & Primary lifecycle with mandatory documentation gates \\
Agile & Adaptive delivery and stakeholder feedback & Visibility and
responsiveness & Can be misused as no-documentation development & Use
values, backlog, reviews, and iteration \\
Scrum & Roles, events, and fixed sprints & Predictable team cadence &
Five days is too short for a meaningful Scrum program & Do not claim
five one-day sprints \\
Kanban & Visualize flow and limit work in progress & Ideal for short
continuous flow & Less prescriptive about planning & Use for daily work
management \\
DevOps & Integrate development, testing, release, and operations &
Repeatable quality and deployment & Requires automation discipline & Use
CI, environment control, release, and monitoring plan \\
\end{longtable}
}

\section{Selected lifecycle
statement}\label{62-selected-lifecycle-statement}

Use the following wording consistently in the report and presentation:

\begin{quote}
UNI-TEL used Rapid Application Development as the primary five-day
lifecycle. The team organized work as daily iterative and incremental
deliveries, managed flow through Kanban, used the imported interface as
a prototype for validation, linked requirements to tests using
V-Model-style traceability, and applied DevOps practices for continuous
integration and controlled release.
\end{quote}

\section{Five-day lifecycle map}\label{63-five-day-lifecycle-map}

\begin{Shaded}
\begin{Highlighting}[]
\NormalTok{Day 1: Initiation, feasibility, baseline audit, scope, backlog}
\NormalTok{   {-}\textgreater{} Day 2: Requirements engineering and high{-}level analysis}
\NormalTok{   {-}\textgreater{} Day 3: Detailed design, test{-}first improvements, CI integration}
\NormalTok{   {-}\textgreater{} Day 4: System integration, verification, deployment rehearsal}
\NormalTok{   {-}\textgreater{} Day 5: Regression, release, evaluation, presentation}
\NormalTok{   {-}\textgreater{} Maintenance: controlled future changes, monitoring, support}
\end{Highlighting}
\end{Shaded}

This sequence displays all classical SDLC phases while remaining honest
about the compressed, iterative execution.

\chapter{Team Organization}\label{7-team-organization}

\section{Primary roles}\label{71-primary-roles}

{\def\LTcaptype{none} % do not increment counter
\begin{longtable}[]{@{}
  >{\raggedright\arraybackslash}p{(\linewidth - 8\tabcolsep) * \real{0.0479}}
  >{\raggedright\arraybackslash}p{(\linewidth - 8\tabcolsep) * \real{0.2335}}
  >{\raggedright\arraybackslash}p{(\linewidth - 8\tabcolsep) * \real{0.2395}}
  >{\raggedright\arraybackslash}p{(\linewidth - 8\tabcolsep) * \real{0.2395}}
  >{\raggedright\arraybackslash}p{(\linewidth - 8\tabcolsep) * \real{0.2395}}@{}}
\toprule\noalign{}
\begin{minipage}[b]{\linewidth}\raggedright
Member
\end{minipage} & \begin{minipage}[b]{\linewidth}\raggedright
Primary role
\end{minipage} & \begin{minipage}[b]{\linewidth}\raggedright
Secondary technical ownership
\end{minipage} & \begin{minipage}[b]{\linewidth}\raggedright
Main documentation ownership
\end{minipage} & \begin{minipage}[b]{\linewidth}\raggedright
Presentation ownership
\end{minipage} \\
\midrule\noalign{}
\endhead
\bottomrule\noalign{}
\endlastfoot
Member 1 & Project Lead and Product Owner & Authentication/security
audit and approved fixes & Charter, scope, PRD, feasibility, backlog &
Problem, planning, feasibility, methodology \\
Member 2 & Business Analyst and UX Lead & Form validation,
accessibility, responsive/UI fixes & SRS, use cases, UI specification &
Requirements, use cases, UI \\
Member 3 & Solution Architect and Data Lead & Supabase schema,
migrations, integrity, database fixes & SDD, architecture, UML/DFD, ERD,
data dictionary & Architecture, database, design \\
Member 4 & QA and Security Test Lead & Test framework, automated tests,
security/system tests & Test plan, cases, defect log, RTM & Testing,
traceability, quality \\
Member 5 & Configuration, DevOps, and Release Lead & CI, environments,
build, deployment, release & SCM plan, deployment, maintenance, release
notes & Git, CI/CD, deployment, maintenance \\
\end{longtable}
}

\section{Fair-contribution rule}\label{72-fair-contribution-rule}

Every member must complete all four categories:

\begin{itemize}
\tightlist
\item
  At least one technical contribution or verified technical analysis.
\item
  At least one authored or co-authored SE artifact.
\item
  At least one pull request reviewed for another member.
\item
  One final presentation segment and participation in the demo or
  questions.
\end{itemize}

Do not use commit count as the only measure of contribution. Use
completed issues, complexity, authored artifacts, tests, reviews,
presentation work, and evidence quality.

\section{RACI matrix}\label{73-raci-matrix}

Key: A = Accountable, R = Responsible, C = Consulted, I = Informed.

{\def\LTcaptype{none} % do not increment counter
\begin{longtable}[]{@{}
  >{\raggedright\arraybackslash}p{(\linewidth - 10\tabcolsep) * \real{0.4776}}
  >{\raggedleft\arraybackslash}p{(\linewidth - 10\tabcolsep) * \real{0.1045}}
  >{\raggedleft\arraybackslash}p{(\linewidth - 10\tabcolsep) * \real{0.1045}}
  >{\raggedleft\arraybackslash}p{(\linewidth - 10\tabcolsep) * \real{0.1045}}
  >{\raggedleft\arraybackslash}p{(\linewidth - 10\tabcolsep) * \real{0.1045}}
  >{\raggedleft\arraybackslash}p{(\linewidth - 10\tabcolsep) * \real{0.1045}}@{}}
\toprule\noalign{}
\begin{minipage}[b]{\linewidth}\raggedright
Activity
\end{minipage} & \begin{minipage}[b]{\linewidth}\raggedleft
M1
\end{minipage} & \begin{minipage}[b]{\linewidth}\raggedleft
M2
\end{minipage} & \begin{minipage}[b]{\linewidth}\raggedleft
M3
\end{minipage} & \begin{minipage}[b]{\linewidth}\raggedleft
M4
\end{minipage} & \begin{minipage}[b]{\linewidth}\raggedleft
M5
\end{minipage} \\
\midrule\noalign{}
\endhead
\bottomrule\noalign{}
\endlastfoot
Scope and charter & A/R & C & C & C & C \\
PRD and backlog & A/R & R & C & C & I \\
SRS and use cases & A & A/R & C & C & I \\
Architecture and database design & C & C & A/R & C & C \\
UI/UX specification & C & A/R & C & C & I \\
Test strategy and RTM & C & C & C & A/R & C \\
Security assessment & R & C & R & A & C \\
Configuration and CI & I & I & C & C & A/R \\
Deployment and release & A & I & C & C & A/R \\
Final report integration & A/R & R & R & R & R \\
Presentation and demonstration & A & R & R & R & R \\
\end{longtable}
}

\section{Review rotation}\label{74-review-rotation}

{\def\LTcaptype{none} % do not increment counter
\begin{longtable}[]{@{}
  >{\raggedright\arraybackslash}p{(\linewidth - 4\tabcolsep) * \real{0.2051}}
  >{\raggedright\arraybackslash}p{(\linewidth - 4\tabcolsep) * \real{0.4103}}
  >{\raggedright\arraybackslash}p{(\linewidth - 4\tabcolsep) * \real{0.3846}}@{}}
\toprule\noalign{}
\begin{minipage}[b]{\linewidth}\raggedright
Author
\end{minipage} & \begin{minipage}[b]{\linewidth}\raggedright
Primary reviewer
\end{minipage} & \begin{minipage}[b]{\linewidth}\raggedright
Backup reviewer
\end{minipage} \\
\midrule\noalign{}
\endhead
\bottomrule\noalign{}
\endlastfoot
Member 1 & Member 2 & Member 4 \\
Member 2 & Member 3 & Member 1 \\
Member 3 & Member 4 & Member 5 \\
Member 4 & Member 5 & Member 3 \\
Member 5 & Member 1 & Member 2 \\
\end{longtable}
}

The primary reviewer checks correctness and traceability. The backup
reviewer is used for high-risk schema, authentication, CI, or release
changes.

\chapter{Project Governance}\label{8-project-governance}

\section{Daily cadence in IST}\label{81-daily-cadence-in-ist}

{\def\LTcaptype{none} % do not increment counter
\begin{longtable}[]{@{}
  >{\raggedright\arraybackslash}p{(\linewidth - 6\tabcolsep) * \real{0.0737}}
  >{\raggedright\arraybackslash}p{(\linewidth - 6\tabcolsep) * \real{0.3368}}
  >{\raggedleft\arraybackslash}p{(\linewidth - 6\tabcolsep) * \real{0.1684}}
  >{\raggedright\arraybackslash}p{(\linewidth - 6\tabcolsep) * \real{0.4211}}@{}}
\toprule\noalign{}
\begin{minipage}[b]{\linewidth}\raggedright
Time
\end{minipage} & \begin{minipage}[b]{\linewidth}\raggedright
Event
\end{minipage} & \begin{minipage}[b]{\linewidth}\raggedleft
Maximum duration
\end{minipage} & \begin{minipage}[b]{\linewidth}\raggedright
Required output
\end{minipage} \\
\midrule\noalign{}
\endhead
\bottomrule\noalign{}
\endlastfoot
09:00 & Daily stand-up & 15 min & Yesterday/today/blockers note \\
09:15 & Planning and task confirmation & 15 min & Updated owners and
board \\
13:00 & Integration checkpoint & 20 min & Early conflicts and dependency
update \\
17:30 & Pull-request and artifact review & 45 min & Review comments and
approvals \\
18:15 & Daily system integration & 30 min & Green main branch or
rollback decision \\
18:45 & Daily review and retrospective & 30 min & Gate result, risk
update, next-day priorities \\
\end{longtable}
}

If the team uses different working hours, preserve the event order and
outputs.

\section{Decision authority}\label{82-decision-authority}

\begin{itemize}
\tightlist
\item
  Product scope and priority: Member 1 after team consultation.
\item
  Requirement wording: Member 2, approved by Member 1.
\item
  Architecture and schema: Member 3, reviewed by Members 4 and 5.
\item
  Release quality and test status: Member 4; Member 1 accepts residual
  academic-project risk.
\item
  CI, deployment, and versioning: Member 5, approved by Member 1.
\item
  Any scope increase after Day 2: formal change request required.
\end{itemize}

\section{Definition of Ready}\label{83-definition-of-ready}

An issue may enter \texttt{In\ Progress} only when it has:

\begin{itemize}
\tightlist
\item
  A clear title and one accountable owner.
\item
  A relevant work-package or requirement ID.
\item
  Testable acceptance criteria.
\item
  Known dependencies and reviewer.
\item
  A size small enough to complete within one day.
\item
  Evidence expectations, such as a document section, test output,
  diagram, screenshot, or deployment URL.
\end{itemize}

\section{Definition of Done}\label{84-definition-of-done}

An issue is Done only when:

\begin{itemize}
\tightlist
\item
  Acceptance criteria are satisfied.
\item
  Requirement and design references are updated where relevant.
\item
  Tests have been written and executed at the appropriate level.
\item
  Type check, lint, automated tests, and production build pass.
\item
  No secrets, personal data, generated junk, or unrelated changes are
  committed.
\item
  A teammate has reviewed the pull request.
\item
  User-facing or operational documentation is updated.
\item
  Evidence is linked to the issue.
\item
  The change is merged and verified on \texttt{main}.
\end{itemize}

\chapter{Repository and Configuration
Management}\label{9-repository-and-configuration-management}

\section{Target repository
structure}\label{91-target-repository-structure}

\begin{Shaded}
\begin{Highlighting}[]
\NormalTok{UNI{-}TEL/}
\NormalTok{|{-}{-} .github/}
\NormalTok{|   |{-}{-} ISSUE\_TEMPLATE/}
\NormalTok{|   |{-}{-} workflows/quality.yml}
\NormalTok{|   \textasciigrave{}{-}{-} pull\_request\_template.md}
\NormalTok{|{-}{-} docs/}
\NormalTok{|   |{-}{-} 01{-}initiation/}
\NormalTok{|   |{-}{-} 02{-}requirements/}
\NormalTok{|   |{-}{-} 03{-}design/}
\NormalTok{|   |{-}{-} 04{-}testing/}
\NormalTok{|   |{-}{-} 05{-}project{-}management/}
\NormalTok{|   |{-}{-} 06{-}deployment{-}maintenance/}
\NormalTok{|   |{-}{-} 07{-}user{-}guide/}
\NormalTok{|   |{-}{-} 08{-}final{-}report/}
\NormalTok{|   |{-}{-} 09{-}presentation/}
\NormalTok{|   |{-}{-} evidence/day{-}1/ ... day{-}5/}
\NormalTok{|   \textasciigrave{}{-}{-} planning/}
\NormalTok{|{-}{-} src/}
\NormalTok{|{-}{-} supabase/}
\NormalTok{|{-}{-} tests/}
\NormalTok{|   |{-}{-} unit/}
\NormalTok{|   |{-}{-} integration/}
\NormalTok{|   \textasciigrave{}{-}{-} fixtures/}
\NormalTok{|{-}{-} output/pdf/}
\NormalTok{|{-}{-} .env.example}
\NormalTok{|{-}{-} .gitignore}
\NormalTok{|{-}{-} CONTRIBUTING.md}
\NormalTok{|{-}{-} LICENSE}
\NormalTok{|{-}{-} README.md}
\NormalTok{\textasciigrave{}{-}{-} package.json}
\end{Highlighting}
\end{Shaded}

\section{Branch naming}\label{92-branch-naming}

\begin{Shaded}
\begin{Highlighting}[]
\NormalTok{docs/WP1{-}charter{-}prd}
\NormalTok{docs/WP2{-}srs{-}use{-}cases}
\NormalTok{design/WP3{-}architecture{-}database}
\NormalTok{test/WP5{-}grade{-}validation}
\NormalTok{fix/WP4{-}schema{-}consistency}
\NormalTok{feat/WP4{-}accessibility{-}validation}
\NormalTok{ci/WP6{-}quality{-}pipeline}
\NormalTok{release/v1.0.0{-}se{-}class}
\end{Highlighting}
\end{Shaded}

\section{Commit conventions}\label{93-commit-conventions}

Use Conventional Commit style and one coherent purpose per commit:

\begin{Shaded}
\begin{Highlighting}[]
\NormalTok{docs(charter): define scope stakeholders and constraints}
\NormalTok{docs(srs): specify FR{-}01 through FR{-}12}
\NormalTok{design(database): add verified ERD and data dictionary}
\NormalTok{test(grades): cover SGPA and weighted marks boundaries}
\NormalTok{fix(validation): reject obtained marks above total marks}
\NormalTok{fix(schema): align schema documentation with migrations}
\NormalTok{feat(a11y): label marks and attendance form controls}
\NormalTok{ci: verify types lint tests and production build}
\NormalTok{docs(rtm): map approved requirements to tests and evidence}
\NormalTok{release: prepare v1.0.0 SE class submission}
\end{Highlighting}
\end{Shaded}

Aim for coherent commits, not a quota. Normal merge commits should
preserve individual authors when contribution history is assessed.

\section{Pull-request requirements}\label{94-pull-request-requirements}

Every PR must contain:

\begin{itemize}
\tightlist
\item
  Linked issue using \texttt{Closes\ \#number} when appropriate.
\item
  Work-package and requirement IDs.
\item
  Concise summary and explicit out-of-scope statement.
\item
  Verification commands and actual results.
\item
  Screenshots for visible UI changes.
\item
  Migration and rollback notes for database changes.
\item
  Documentation and RTM impact.
\item
  At least one approving reviewer.
\end{itemize}

\section{Kanban board}\label{95-kanban-board}

Use these columns and WIP limits:

{\def\LTcaptype{none} % do not increment counter
\begin{longtable}[]{@{}
  >{\raggedright\arraybackslash}p{(\linewidth - 4\tabcolsep) * \real{0.2459}}
  >{\raggedright\arraybackslash}p{(\linewidth - 4\tabcolsep) * \real{0.5574}}
  >{\raggedleft\arraybackslash}p{(\linewidth - 4\tabcolsep) * \real{0.1967}}@{}}
\toprule\noalign{}
\begin{minipage}[b]{\linewidth}\raggedright
Column
\end{minipage} & \begin{minipage}[b]{\linewidth}\raggedright
Meaning
\end{minipage} & \begin{minipage}[b]{\linewidth}\raggedleft
WIP limit
\end{minipage} \\
\midrule\noalign{}
\endhead
\bottomrule\noalign{}
\endlastfoot
Product Backlog & Approved but not scheduled & None \\
Selected & Ready for current day & 10 \\
In Progress & Actively owned & 1 per member \\
In Review & PR or document awaiting review & 5 \\
Testing & Integrated verification & 5 \\
Done & Meets Definition of Done & None \\
Blocked & Waiting for dependency or decision & None \\
\end{longtable}
}

\section{Version and change
control}\label{96-version-and-change-control}

\begin{itemize}
\tightlist
\item
  Baseline: \texttt{0.1.0-baseline} after imported prototype
  verification.
\item
  Candidate: \texttt{1.0.0-rc.1} after Day 4 integration.
\item
  Final: \texttt{v1.0.0-se-class} after Day 5 release gate.
\item
  Documents carry version, author, reviewer, approval date, and change
  history.
\item
  Schema changes use new forward migrations; do not silently edit an
  already-applied migration.
\item
  Secrets exist only in local or hosted environment settings, never in
  Git.
\end{itemize}

\chapter{Work Breakdown Structure}\label{10-work-breakdown-structure}

{\def\LTcaptype{none} % do not increment counter
\begin{longtable}[]{@{}
  >{\raggedright\arraybackslash}p{(\linewidth - 10\tabcolsep) * \real{0.0625}}
  >{\raggedright\arraybackslash}p{(\linewidth - 10\tabcolsep) * \real{0.3036}}
  >{\raggedright\arraybackslash}p{(\linewidth - 10\tabcolsep) * \real{0.3571}}
  >{\raggedright\arraybackslash}p{(\linewidth - 10\tabcolsep) * \real{0.0714}}
  >{\raggedright\arraybackslash}p{(\linewidth - 10\tabcolsep) * \real{0.1429}}
  >{\raggedright\arraybackslash}p{(\linewidth - 10\tabcolsep) * \real{0.0625}}@{}}
\toprule\noalign{}
\begin{minipage}[b]{\linewidth}\raggedright
WP
\end{minipage} & \begin{minipage}[b]{\linewidth}\raggedright
Work package
\end{minipage} & \begin{minipage}[b]{\linewidth}\raggedright
Main outcome
\end{minipage} & \begin{minipage}[b]{\linewidth}\raggedright
Owner
\end{minipage} & \begin{minipage}[b]{\linewidth}\raggedright
Primary reviewer
\end{minipage} & \begin{minipage}[b]{\linewidth}\raggedright
Due
\end{minipage} \\
\midrule\noalign{}
\endhead
\bottomrule\noalign{}
\endlastfoot
WP0 & Repository and baseline governance & Clean baseline, board,
templates, contribution rules & M5 & M1 & Day 1 \\
WP1 & Initiation and product planning & Charter, feasibility, PRD,
scope, backlog & M1 & M2 & Day 1 \\
WP2 & Requirements engineering & SRS, actors, use cases, NFRs,
acceptance criteria & M2 & M3 & Day 2 \\
WP3 & Analysis and design & SDD, UML, DFD, ERD, data dictionary,
UI/security design & M3 & M4 & Day 3 \\
WP4 & Targeted product hardening & Verified fixes for security,
validation, schema, UX & M1/M2/M3 & M4 & Day 4 \\
WP5 & Verification and validation & Automated/manual tests, defect log,
RTM, report & M4 & M5 & Day 4 \\
WP6 & CI/CD and release engineering & Quality workflow, environments,
deploy, release/rollback & M5 & M1 & Day 5 \\
WP7 & Management controls & Gantt, estimates, risks, changes, meetings,
metrics & M1 & M5 & Daily \\
WP8 & Final report and presentation & Integrated report, slides, demo,
viva, release package & All & M1 & Day 5 \\
\end{longtable}
}

\chapter{Prioritized Product and SE
Backlog}\label{11-prioritized-product-and-se-backlog}

Use MoSCoW priority. Story points are relative planning values, not
elapsed hours.

{\def\LTcaptype{none} % do not increment counter
\begin{longtable}[]{@{}
  >{\raggedright\arraybackslash}p{(\linewidth - 14\tabcolsep) * \real{0.0680}}
  >{\raggedright\arraybackslash}p{(\linewidth - 14\tabcolsep) * \real{0.3883}}
  >{\raggedright\arraybackslash}p{(\linewidth - 14\tabcolsep) * \real{0.1262}}
  >{\raggedleft\arraybackslash}p{(\linewidth - 14\tabcolsep) * \real{0.0680}}
  >{\raggedright\arraybackslash}p{(\linewidth - 14\tabcolsep) * \real{0.0680}}
  >{\raggedright\arraybackslash}p{(\linewidth - 14\tabcolsep) * \real{0.0777}}
  >{\raggedright\arraybackslash}p{(\linewidth - 14\tabcolsep) * \real{0.0680}}
  >{\raggedright\arraybackslash}p{(\linewidth - 14\tabcolsep) * \real{0.1359}}@{}}
\toprule\noalign{}
\begin{minipage}[b]{\linewidth}\raggedright
ID
\end{minipage} & \begin{minipage}[b]{\linewidth}\raggedright
Backlog item
\end{minipage} & \begin{minipage}[b]{\linewidth}\raggedright
Priority
\end{minipage} & \begin{minipage}[b]{\linewidth}\raggedleft
Points
\end{minipage} & \begin{minipage}[b]{\linewidth}\raggedright
Owner
\end{minipage} & \begin{minipage}[b]{\linewidth}\raggedright
Reviewer
\end{minipage} & \begin{minipage}[b]{\linewidth}\raggedright
Target
\end{minipage} & \begin{minipage}[b]{\linewidth}\raggedright
Dependency
\end{minipage} \\
\midrule\noalign{}
\endhead
\bottomrule\noalign{}
\endlastfoot
SE-001 & Import and verify clean baseline & Must & 3 & M5 & M1 & D1 &
None \\
SE-002 & Create charter, scope, objectives, stakeholder register & Must
& 5 & M1 & M2 & D1 & SE-001 \\
SE-003 & Complete five-part feasibility study & Must & 3 & M1 & M3 & D1
& SE-002 \\
SE-004 & Compare lifecycle models and record RAD decision & Must & 3 &
M1 & M4 & D1 & SE-002 \\
SE-005 & Audit features, code, schema, tests, docs, security, deployment
& Must & 8 & All & M1 & D1 & SE-001 \\
SE-006 & Configure board, templates, contribution and review policy &
Must & 5 & M5 & M1 & D1 & SE-001 \\
SE-007 & Write PRD and prioritized product backlog & Must & 8 & M1 & M2
& D2 & SE-005 \\
SE-008 & Write SRS with FR/NFR identifiers and business rules & Must &
13 & M2 & M3 & D2 & SE-007 \\
SE-009 & Define actors, use cases, acceptance criteria, and UI flows &
Must & 8 & M2 & M1 & D2 & SE-008 \\
SE-010 & Create context, use-case, DFD 0/1, activity diagrams & Must & 8
& M3 & M2 & D3 & SE-008 \\
SE-011 & Create architecture, component, sequence, deployment diagrams &
Must & 8 & M3 & M4 & D3 & SE-008 \\
SE-012 & Verify ERD, schema, constraints, RLS, and data dictionary &
Must & 13 & M3 & M4 & D3 & SE-005 \\
SE-013 & Create test strategy, test data, cases, and RTM skeleton & Must
& 8 & M4 & M5 & D2 & SE-008 \\
SE-014 & Add automated test framework and calculation/validation tests &
Must & 13 & M4 & M5 & D3 & SE-013 \\
SE-015 & Add type-check/test scripts and GitHub Actions quality gate &
Must & 8 & M5 & M4 & D3 & SE-014 \\
SE-016 & Correct verified security/authentication issues & Must if found
& 8 & M1 & M4 & D4 & SE-005 \\
SE-017 & Correct verified validation/accessibility/responsive issues &
Must if found & 8 & M2 & M4 & D4 & SE-005 \\
SE-018 & Correct verified schema/integrity/documentation issues & Must
if found & 8 & M3 & M4 & D4 & SE-012 \\
SE-019 & Execute unit, integration, system, security, UX tests & Must &
13 & M4 & M1 & D4 & SE-014 to 018 \\
SE-020 & Manage defect triage, retest, regression, and RTM completion &
Must & 8 & M4 & M5 & D5 & SE-019 \\
SE-021 & Document configuration, deployment, rollback, maintenance &
Must & 8 & M5 & M1 & D4 & SE-015 \\
SE-022 & Deploy release candidate and preserve evidence & Must & 5 & M5
& M1 & D4 & SE-015, SE-021 \\
SE-023 & Create user manual and installation guide & Must & 5 & M2/M5 &
M3 & D5 & SE-022 \\
SE-024 & Integrate final report and appendices & Must & 13 & All & M1 &
D5 & SE-002 to 023 \\
SE-025 & Prepare slides, demo, viva, and contribution report & Must & 8
& All & M1 & D5 & SE-024 \\
SE-026 & Tag release and archive final submission & Must & 3 & M5 & M1 &
D5 & SE-020, SE-025 \\
\end{longtable}
}

The team should not begin Could-have features until every Must item
required for the release gate is complete.

\chapter{Five-Day Execution
Schedule}\label{12-five-day-execution-schedule}

\section{Day 1 - Friday, 14 August
2026}\label{day-1---friday-14-august-2026}

\textbf{Theme:} Initiation, transparent baseline, feasibility, scope,
and governance.

\textbf{Team objectives:} Establish the official repository, understand
the imported system, decide what will and will not be delivered, assign
ownership, and create an auditable plan.

{\def\LTcaptype{none} % do not increment counter
\begin{longtable}[]{@{}
  >{\raggedright\arraybackslash}p{(\linewidth - 4\tabcolsep) * \real{0.0805}}
  >{\raggedright\arraybackslash}p{(\linewidth - 4\tabcolsep) * \real{0.4598}}
  >{\raggedright\arraybackslash}p{(\linewidth - 4\tabcolsep) * \real{0.4598}}@{}}
\toprule\noalign{}
\begin{minipage}[b]{\linewidth}\raggedright
Member
\end{minipage} & \begin{minipage}[b]{\linewidth}\raggedright
Assigned work
\end{minipage} & \begin{minipage}[b]{\linewidth}\raggedright
Required evidence
\end{minipage} \\
\midrule\noalign{}
\endhead
\bottomrule\noalign{}
\endlastfoot
M1 & Draft charter, problem, vision, objectives, stakeholders, scope,
assumptions, constraints; lead methodology decision & Charter PR,
decision record ADR-001, issue links \\
M2 & Audit current user flows, forms, pages, responsive behavior,
accessibility; identify actors and preliminary use cases & UI audit with
screenshots and issue list \\
M3 & Audit architecture, services, migrations, entities, constraints,
RLS; flag inconsistencies & Architecture/database audit and verified
schema notes \\
M4 & Audit existing tests, calculation risks, validation, security
scenarios, and quality claims & Test-gap report and initial risk/test
inventory \\
M5 & Create clean baseline, \texttt{.gitignore}, \texttt{.env.example},
templates, board, milestone, branch policy, build audit & Baseline tag,
repository settings evidence, build log \\
\end{longtable}
}

\textbf{Whole-team activities:}

\begin{itemize}
\tightlist
\item
  Demonstrate the baseline from login through analytics/export.
\item
  Agree on scope freeze and selected methodology.
\item
  Convert audit findings into prioritized issues.
\item
  Assign one owner and reviewer to every Day 2 issue.
\item
  Establish document naming, ID conventions, and evidence locations.
\end{itemize}

\textbf{Required outputs:}

\begin{itemize}
\tightlist
\item
  Project Charter v1.0.
\item
  Stakeholder Register.
\item
  Technical, economic, operational, schedule, and legal/security
  feasibility study.
\item
  Methodology comparison and ADR-001 selecting RAD.
\item
  Scope statement and feature freeze.
\item
  Initial product/SE backlog and five-day Gantt chart.
\item
  Risk Register v0.1.
\item
  Baseline audit reports.
\item
  Repository contribution, review, and configuration rules.
\end{itemize}

\textbf{Suggested commits:}

\begin{Shaded}
\begin{Highlighting}[]
\NormalTok{chore: import UNI{-}TEL baseline prototype}
\NormalTok{chore(repo): add contribution and review controls}
\NormalTok{docs(charter): define project scope and stakeholders}
\NormalTok{docs(feasibility): assess five{-}day prototype viability}
\NormalTok{docs(methodology): select RAD with traceability practices}
\NormalTok{docs(audit): record baseline UI architecture test and release gaps}
\end{Highlighting}
\end{Shaded}

\textbf{Day 1 exit gate:} Baseline builds; no secrets are tracked; scope
is frozen; methodology is approved; responsibilities and reviewers are
assigned; Day 2 issues satisfy Definition of Ready.

\section{Day 2 - Saturday, 15 August
2026}\label{day-2---saturday-15-august-2026}

\textbf{Theme:} Requirements engineering, analysis, acceptance, and test
planning.

\textbf{Team objectives:} Convert the product vision and baseline
behavior into approved, numbered, measurable requirements that can drive
design and testing.

{\def\LTcaptype{none} % do not increment counter
\begin{longtable}[]{@{}
  >{\raggedright\arraybackslash}p{(\linewidth - 4\tabcolsep) * \real{0.0805}}
  >{\raggedright\arraybackslash}p{(\linewidth - 4\tabcolsep) * \real{0.4598}}
  >{\raggedright\arraybackslash}p{(\linewidth - 4\tabcolsep) * \real{0.4598}}@{}}
\toprule\noalign{}
\begin{minipage}[b]{\linewidth}\raggedright
Member
\end{minipage} & \begin{minipage}[b]{\linewidth}\raggedright
Assigned work
\end{minipage} & \begin{minipage}[b]{\linewidth}\raggedright
Required evidence
\end{minipage} \\
\midrule\noalign{}
\endhead
\bottomrule\noalign{}
\endlastfoot
M1 & Complete PRD, personas/user goals, business value, MoSCoW
priorities, release criteria & PRD PR and backlog mapping \\
M2 & Complete SRS, actors, FRs, NFRs, business rules, detailed use
cases, acceptance criteria & SRS PR with review comments resolved \\
M3 & Create system context and initial logical data/process models;
verify requirement feasibility against architecture & Context/DFD drafts
and feasibility comments \\
M4 & Create test strategy, test levels/types, entry/exit criteria, test
environment, test data, RTM skeleton & Test Plan v0.1 and initial RTM \\
M5 & Create SCM plan, CI design, environment inventory, deployment
architecture draft, backup/rollback plan & SCM/DevOps documents and
workflow issue \\
\end{longtable}
}

\textbf{Formal requirements review:} At 17:30, read each Must
requirement. Check that it is necessary, singular, unambiguous,
feasible, testable, uniquely identified, prioritized, and traceable.

\textbf{Required outputs:}

\begin{itemize}
\tightlist
\item
  PRD v1.0.
\item
  SRS v1.0.
\item
  Actor and use-case catalogue.
\item
  Use-case diagram and detailed descriptions for critical workflows.
\item
  Functional and non-functional requirement catalogues.
\item
  Business rules and data-validation rules.
\item
  MoSCoW priority matrix.
\item
  Acceptance criteria.
\item
  Test Plan v0.1 and RTM v0.1.
\item
  Requirements Review Minutes and approved baseline.
\end{itemize}

\textbf{Suggested commits:}

\begin{Shaded}
\begin{Highlighting}[]
\NormalTok{docs(prd): define users value scope and release outcomes}
\NormalTok{docs(srs): specify functional and nonfunctional requirements}
\NormalTok{docs(use{-}cases): model core student workflows}
\NormalTok{test(plan): define verification strategy and RTM structure}
\NormalTok{docs(scm): define branches environments builds and releases}
\end{Highlighting}
\end{Shaded}

\textbf{Day 2 exit gate:} All Must requirements are approved and
testable; requirement IDs are frozen; every Must requirement has
acceptance criteria and planned tests; open ambiguity is recorded as a
decision or change request.

\section{Day 3 - Sunday, 16 August
2026}\label{day-3---sunday-16-august-2026}

\textbf{Theme:} Detailed design, test-first product hardening,
automation, and design review.

\textbf{Team objectives:} Produce consistent architecture and data
designs, add the missing automated quality foundation, and implement
only verified high-value corrections.

{\def\LTcaptype{none} % do not increment counter
\begin{longtable}[]{@{}
  >{\raggedright\arraybackslash}p{(\linewidth - 4\tabcolsep) * \real{0.0805}}
  >{\raggedright\arraybackslash}p{(\linewidth - 4\tabcolsep) * \real{0.4598}}
  >{\raggedright\arraybackslash}p{(\linewidth - 4\tabcolsep) * \real{0.4598}}@{}}
\toprule\noalign{}
\begin{minipage}[b]{\linewidth}\raggedright
Member
\end{minipage} & \begin{minipage}[b]{\linewidth}\raggedright
Assigned work
\end{minipage} & \begin{minipage}[b]{\linewidth}\raggedright
Required evidence
\end{minipage} \\
\midrule\noalign{}
\endhead
\bottomrule\noalign{}
\endlastfoot
M1 & Verify authentication/session/authorization design; write failing
security/route scenarios with M4; implement approved fix if a failure
exists & Issue, test evidence, reviewed PR \\
M2 & Finalize UI specification and activity flows; write
validation/accessibility scenarios; implement approved issue using
test-first steps & Before/after evidence and reviewed PR \\
M3 & Complete SDD, ERD, data dictionary, constraints/RLS, UML/DFD; fix
only verified schema or schema-documentation defects & Design package
and migration/verification evidence \\
M4 & Configure Vitest/RTL; add unit tests for grade, attendance, marks,
semester normalization, and validation services & Failing-then-passing
test commits and coverage output \\
M5 & Add type-check/test scripts and GitHub Actions workflow; verify
clean install and production build & Green workflow link/log and
failure-handling notes \\
\end{longtable}
}

\textbf{Required design views:} context, use-case, DFD Level 0, DFD
Level 1, activity, sequence, component, ER, deployment, and high-level
architecture. Every diagram must contain a title, legend where needed,
consistent names, and a paragraph explaining how it relates to
requirements.

\textbf{Test-first change cycle:}

\begin{enumerate}
\def\labelenumi{\arabic{enumi}.}
\tightlist
\item
  Link an issue to a requirement or verified audit finding.
\item
  Add a test or reproducible failure procedure.
\item
  Demonstrate that the check fails for the intended reason.
\item
  Apply the smallest correction.
\item
  Demonstrate that the check passes and no regression appears.
\item
  Update documentation and RTM.
\item
  Obtain peer review before merge.
\end{enumerate}

\textbf{Required outputs:}

\begin{itemize}
\tightlist
\item
  Software Design Document v1.0.
\item
  Complete diagram set and data dictionary.
\item
  UI, validation, error-handling, and security design sections.
\item
  Automated test framework and initial critical tests.
\item
  CI quality workflow.
\item
  Verified technical correction PRs, if failures were found.
\item
  Architecture and Design Review Minutes.
\end{itemize}

\textbf{Suggested commits:}

\begin{Shaded}
\begin{Highlighting}[]
\NormalTok{design(architecture): document component and deployment views}
\NormalTok{design(database): add verified ERD constraints RLS and dictionary}
\NormalTok{test(setup): configure Vitest and React Testing Library}
\NormalTok{test(calculations): cover SGPA attendance and marks boundaries}
\NormalTok{test(import): cover semester normalization regressions}
\NormalTok{ci: run type check lint tests and build on pull requests}
\NormalTok{fix(...): correct the verified behavior described by the linked issue}
\end{Highlighting}
\end{Shaded}

\textbf{Day 3 exit gate:} Design names agree with SRS and code; critical
automated tests pass; CI is green; every correction has failure evidence
and review; design changes are reflected in the RTM.

\section{Day 4 - Monday, 17 August
2026}\label{day-4---monday-17-august-2026}

\textbf{Theme:} Integration, system verification, defect resolution, and
release rehearsal.

\textbf{Team objectives:} Verify complete user journeys and
non-functional targets, close release-blocking defects, complete
deployment documentation, and create the release candidate.

{\def\LTcaptype{none} % do not increment counter
\begin{longtable}[]{@{}
  >{\raggedright\arraybackslash}p{(\linewidth - 4\tabcolsep) * \real{0.0805}}
  >{\raggedright\arraybackslash}p{(\linewidth - 4\tabcolsep) * \real{0.4598}}
  >{\raggedright\arraybackslash}p{(\linewidth - 4\tabcolsep) * \real{0.4598}}@{}}
\toprule\noalign{}
\begin{minipage}[b]{\linewidth}\raggedright
Member
\end{minipage} & \begin{minipage}[b]{\linewidth}\raggedright
Assigned work
\end{minipage} & \begin{minipage}[b]{\linewidth}\raggedright
Required evidence
\end{minipage} \\
\midrule\noalign{}
\endhead
\bottomrule\noalign{}
\endlastfoot
M1 & Execute acceptance tests for scope and critical journeys; triage
defects; control scope/change requests & UAT record, decisions, updated
risk register \\
M2 & Execute usability, accessibility, responsive, browser,
form-validation, and user-manual walkthrough tests & Device/browser
matrix, screenshots, issue results \\
M3 & Execute database CRUD, constraint, cascade, calculation, migration,
and RLS integration tests & Query/test log using synthetic users and
data \\
M4 & Coordinate unit, integration, system, regression, security,
negative, recovery, and exploratory testing; maintain defect log and RTM
& Test execution report, defect dashboard, RTM v0.9 \\
M5 & Run clean build, environment and deployment rehearsal; deploy
\texttt{1.0.0-rc.1}; test rollback and document operations & Deployment
URL/evidence, build log, rollback record \\
\end{longtable}
}

\textbf{Core end-to-end scenarios:}

\begin{enumerate}
\def\labelenumi{\arabic{enumi}.}
\tightlist
\item
  Register/login and reach protected dashboard.
\item
  Create semester and add subjects/credits.
\item
  Enter valid and invalid attendance data.
\item
  Enter valid and invalid marks/weightage.
\item
  Verify calculated percentages, SGPA, and CGPA.
\item
  View analytics and empty/error/loading states.
\item
  Import valid data and reject malformed or unsafe data.
\item
  Export supported report formats.
\item
  Verify logout and protected-route behavior.
\item
  Verify User A cannot access User B\textquotesingle s records.
\end{enumerate}

\textbf{Required outputs:}

\begin{itemize}
\tightlist
\item
  Executed test-case register with actual results.
\item
  Automated test and CI reports.
\item
  Defect Log with severity, priority, owner, resolution, retest.
\item
  Security, usability, compatibility, and data-integrity evidence.
\item
  RTM v0.9.
\item
  Deployment Guide, rollback plan, maintenance plan, and operations
  checklist.
\item
  Release Candidate \texttt{1.0.0-rc.1} and rehearsal record.
\end{itemize}

\textbf{Suggested commits:}

\begin{Shaded}
\begin{Highlighting}[]
\NormalTok{test(system): record core workflow verification}
\NormalTok{test(security): verify protected routes and user isolation}
\NormalTok{test(database): verify constraints cascades and RLS}
\NormalTok{fix(...): resolve release{-}blocking defect DEF{-}xxx}
\NormalTok{docs(deploy): add environment deployment and rollback runbook}
\NormalTok{docs(rtm): link requirements design implementation and results}
\NormalTok{release: prepare 1.0.0{-}rc.1}
\end{Highlighting}
\end{Shaded}

\textbf{Day 4 exit gate:} Release candidate builds and deploys; all Must
requirements have test results; no unresolved Critical defects; High
defects are fixed or explicitly risk-accepted; rollback is documented;
RTM has no unexplained Must-requirement gaps.

\section{Day 5 - Tuesday, 18 August
2026}\label{day-5---tuesday-18-august-2026}

\textbf{Theme:} Final regression, controlled release, reporting,
demonstration, and evaluation.

\textbf{Team objectives:} Freeze the product, verify the release,
integrate academic evidence, rehearse presentation/demo, and archive a
reproducible submission.

{\def\LTcaptype{none} % do not increment counter
\begin{longtable}[]{@{}
  >{\raggedright\arraybackslash}p{(\linewidth - 4\tabcolsep) * \real{0.0805}}
  >{\raggedright\arraybackslash}p{(\linewidth - 4\tabcolsep) * \real{0.4598}}
  >{\raggedright\arraybackslash}p{(\linewidth - 4\tabcolsep) * \real{0.4598}}@{}}
\toprule\noalign{}
\begin{minipage}[b]{\linewidth}\raggedright
Member
\end{minipage} & \begin{minipage}[b]{\linewidth}\raggedright
Assigned work
\end{minipage} & \begin{minipage}[b]{\linewidth}\raggedright
Required evidence
\end{minipage} \\
\midrule\noalign{}
\endhead
\bottomrule\noalign{}
\endlastfoot
M1 & Chair release gate; integrate executive report sections; finalize
contribution and lessons-learned records & Approval minutes, report,
contribution matrix \\
M2 & Finalize user manual, screenshots, requirements narrative, slides,
and user-facing demo data & Reviewed guide and slide segment \\
M3 & Verify all design diagrams/data dictionary against final release;
finalize architecture slides and viva notes & Design consistency
checklist \\
M4 & Run final regression and smoke tests; finalize Test Report, Defect
Summary, RTM, and quality metrics & Signed test summary and green result
links \\
M5 & Verify clean clone/build/deploy, tag release, generate release
notes, archive checksums and final package & \texttt{v1.0.0-se-class},
release page, archive manifest \\
\end{longtable}
}

\textbf{Release sequence:}

\begin{enumerate}
\def\labelenumi{\arabic{enumi}.}
\tightlist
\item
  Freeze features and accept documentation/test fixes only.
\item
  Pull a clean copy and install from the lock file.
\item
  Run type checking, lint, tests, and production build.
\item
  Run smoke tests against the release candidate.
\item
  Confirm secrets and personal data are absent.
\item
  Confirm documents, diagrams, RTM, test report, user guide, and slides
  are current.
\item
  Record professor-ready demo data and reset procedure.
\item
  Approve the release in signed meeting minutes.
\item
  Merge release PR and tag \texttt{v1.0.0-se-class}.
\item
  Rehearse the timed presentation, live demo, backup demo, and viva
  transitions.
\end{enumerate}

\textbf{Required outputs:}

\begin{itemize}
\tightlist
\item
  Final Project Report and appendices.
\item
  PRD, SRS, SDD, Test Report, RTM, User Manual, Deployment and
  Maintenance documents.
\item
  Final diagrams in editable and exported formats.
\item
  Presentation slides, speaker allocation, demo script, and viva
  question bank.
\item
  Contribution report based on issues, commits, PRs, reviews, artifacts,
  and speaking roles.
\item
  Final release, release notes, checksum/archive manifest, and
  submission checklist.
\end{itemize}

\textbf{Suggested commits:}

\begin{Shaded}
\begin{Highlighting}[]
\NormalTok{docs(report): integrate final SE project report}
\NormalTok{docs(user{-}guide): finalize guided UNI{-}TEL workflows}
\NormalTok{test(regression): record final release verification}
\NormalTok{docs(presentation): add final slides demo and viva guide}
\NormalTok{docs(contributions): summarize verified team work}
\NormalTok{release: publish v1.0.0{-}se{-}class}
\end{Highlighting}
\end{Shaded}

\textbf{Day 5 exit gate:} Every final checklist item is signed; CI is
green; clean build and deployment are verified; RTM is complete; release
is tagged; all members can explain both theory and their evidence.

\chapter{Required Software Engineering Artifact
Set}\label{13-required-software-engineering-artifact-set}

Each artifact must have a title, version, authors, reviewers, approval,
revision history, table of contents where appropriate, and links to
related requirements or evidence.

{\def\LTcaptype{none} % do not increment counter
\begin{longtable}[]{@{}
  >{\raggedleft\arraybackslash}p{(\linewidth - 8\tabcolsep) * \real{0.0787}}
  >{\raggedright\arraybackslash}p{(\linewidth - 8\tabcolsep) * \real{0.3146}}
  >{\raggedright\arraybackslash}p{(\linewidth - 8\tabcolsep) * \real{0.4494}}
  >{\raggedright\arraybackslash}p{(\linewidth - 8\tabcolsep) * \real{0.0787}}
  >{\raggedright\arraybackslash}p{(\linewidth - 8\tabcolsep) * \real{0.0787}}@{}}
\toprule\noalign{}
\begin{minipage}[b]{\linewidth}\raggedleft
No.
\end{minipage} & \begin{minipage}[b]{\linewidth}\raggedright
Artifact
\end{minipage} & \begin{minipage}[b]{\linewidth}\raggedright
Minimum content
\end{minipage} & \begin{minipage}[b]{\linewidth}\raggedright
Owner
\end{minipage} & \begin{minipage}[b]{\linewidth}\raggedright
Due
\end{minipage} \\
\midrule\noalign{}
\endhead
\bottomrule\noalign{}
\endlastfoot
01 & Project Proposal and Charter & Problem, vision, objectives,
stakeholders, scope, constraints, success, approvals & M1 & D1 \\
02 & Feasibility Study & Technical, economic, operational, schedule,
legal/security conclusions & M1 & D1 \\
03 & Methodology Study & Lifecycle theory comparison, selection
criteria, ADR-001 & M1 & D1 \\
04 & Project Management Plan & WBS, schedule, effort, roles, RACI,
communications, quality, procurement/tools & M1/M5 & D1 \\
05 & PRD & Users, needs, value, features, exclusions, priorities,
metrics, release criteria & M1 & D2 \\
06 & SRS & Introduction, overall description, FRs, NFRs, interfaces,
rules, constraints, acceptance & M2 & D2 \\
07 & Use-Case Specification & Actors, diagram, pre/postconditions,
main/alternate/exception flows & M2 & D2 \\
08 & Software Design Document & Architecture, modules, data flows,
interfaces, security, errors, decisions & M3 & D3 \\
09 & Diagram Set & Context, use case, DFD 0/1, activity, sequence,
component, ERD, deployment & M3 & D3 \\
10 & Database Design & Tables, keys, constraints, indexes, RLS,
dictionary, migration strategy & M3 & D3 \\
11 & UI/UX Specification & Navigation, screen inventory, validation,
states, accessibility, responsive rules & M2 & D3 \\
12 & Test Plan & Scope, levels/types, environment, data, entry/exit,
roles, reporting & M4 & D2 \\
13 & Test Cases and Results & Preconditions, inputs, steps,
expected/actual, status, evidence & M4 & D4 \\
14 & Defect Log and Test Report & Severity, lifecycle, resolution,
retest, metrics, residual risk & M4 & D5 \\
15 & RTM & Requirement to use case, design, implementation, test,
result, evidence & M4 & D5 \\
16 & Risk Management Plan & Identification, scoring, response, owner,
trigger, residual risk & M1/M4 & Daily \\
17 & SCM and Change Plan & Baselines, Git flow, reviews, versions,
builds, changes, audits & M5 & D2 \\
18 & CI/CD and Deployment Guide & Pipeline, variables, build, deploy,
smoke, rollback, recovery & M5 & D4 \\
19 & Maintenance Plan & Corrective, adaptive, perfective, preventive
work; support and monitoring & M5 & D5 \\
20 & User Manual & Setup/login, workflows, validation, troubleshooting,
screenshots & M2 & D5 \\
21 & Final Project Report & Integrated SE narrative, evidence, results,
limitations, learning, references & All & D5 \\
22 & Presentation and Demo & Theory-to-evidence slides, timed script,
backup evidence, viva preparation & All & D5 \\
23 & Team Evidence Pack & Minutes, decisions, board, commits, PRs,
reviews, daily logs, contributions & M1 & D5 \\
\end{longtable}
}

\chapter{Requirements Engineering
Standard}\label{14-requirements-engineering-standard}

\section{ID scheme}\label{141-id-scheme}

\begin{Shaded}
\begin{Highlighting}[]
\NormalTok{BUS{-}xx   Business objective}
\NormalTok{USR{-}xx   User need}
\NormalTok{FR{-}xx    Functional requirement}
\NormalTok{NFR{-}xx   Non{-}functional requirement}
\NormalTok{BR{-}xx    Business rule}
\NormalTok{UC{-}xx    Use case}
\NormalTok{AC{-}xx    Acceptance criterion}
\NormalTok{TC{-}xx    Test case}
\NormalTok{DEF{-}xx   Defect}
\NormalTok{RISK{-}xx  Risk}
\NormalTok{CR{-}xx    Change request}
\NormalTok{ADR{-}xxx  Architecture/project decision record}
\end{Highlighting}
\end{Shaded}

\section{Seed functional
requirements}\label{142-seed-functional-requirements}

These are starting requirements. Member 2 must verify them against the
baseline and scope before the Day 2 review.

{\def\LTcaptype{none} % do not increment counter
\begin{longtable}[]{@{}
  >{\raggedright\arraybackslash}p{(\linewidth - 4\tabcolsep) * \real{0.1273}}
  >{\raggedright\arraybackslash}p{(\linewidth - 4\tabcolsep) * \real{0.7273}}
  >{\raggedright\arraybackslash}p{(\linewidth - 4\tabcolsep) * \real{0.1455}}@{}}
\toprule\noalign{}
\begin{minipage}[b]{\linewidth}\raggedright
ID
\end{minipage} & \begin{minipage}[b]{\linewidth}\raggedright
Requirement
\end{minipage} & \begin{minipage}[b]{\linewidth}\raggedright
Priority
\end{minipage} \\
\midrule\noalign{}
\endhead
\bottomrule\noalign{}
\endlastfoot
FR-01 & The system shall allow a student to register, authenticate, end
a session, and recover from authentication errors safely. & Must \\
FR-02 & The system shall permit authenticated users to access only their
own academic records. & Must \\
FR-03 & The student shall create, view, edit, and delete semesters
subject to integrity rules. & Must \\
FR-04 & The student shall manage subjects, credits, and grades within a
selected semester. & Must \\
FR-05 & The student shall record and update subject attendance values,
and the system shall calculate the percentage. & Must \\
FR-06 & The student shall record examinations, totals, obtained marks,
and weightage subject to validation. & Must \\
FR-07 & The system shall calculate and display valid marks percentages
and weighted performance. & Must \\
FR-08 & The system shall calculate and display semester SGPA and
cumulative CGPA from approved grade rules. & Must \\
FR-09 & The system shall display academic summaries, trends,
distributions, and warnings supported by stored data. & Must \\
FR-10 & The system shall import supported academic data and report
invalid input without corrupting existing records. & Should \\
FR-11 & The system shall export supported academic data or reports in
the formats implemented by the release. & Should \\
FR-12 & The system shall present notifications or warnings for relevant
academic conditions implemented by the release. & Should \\
FR-13 & The student shall view and update supported profile and
preference information. & Could \\
\end{longtable}
}

\section{Seed non-functional
requirements}\label{143-seed-non-functional-requirements}

Values are release targets. Results must be measured under a documented
test environment.

{\def\LTcaptype{none} % do not increment counter
\begin{longtable}[]{@{}
  >{\raggedright\arraybackslash}p{(\linewidth - 4\tabcolsep) * \real{0.1094}}
  >{\raggedright\arraybackslash}p{(\linewidth - 4\tabcolsep) * \real{0.2656}}
  >{\raggedright\arraybackslash}p{(\linewidth - 4\tabcolsep) * \real{0.6250}}@{}}
\toprule\noalign{}
\begin{minipage}[b]{\linewidth}\raggedright
ID
\end{minipage} & \begin{minipage}[b]{\linewidth}\raggedright
Quality attribute
\end{minipage} & \begin{minipage}[b]{\linewidth}\raggedright
Measurable requirement
\end{minipage} \\
\midrule\noalign{}
\endhead
\bottomrule\noalign{}
\endlastfoot
NFR-01 & Security & Unauthenticated protected routes redirect to
authentication, and database RLS prevents cross-user CRUD access in
executed two-user tests. \\
NFR-02 & Data integrity & Marks satisfy
\texttt{0\ \textless{}=\ obtained\ \textless{}=\ total}, total is
positive, credits and percentages satisfy documented ranges, and invalid
inputs are rejected. \\
NFR-03 & Reliability & No uncaught application failure occurs during the
ten critical end-to-end scenarios; recoverable errors show a usable
message. \\
NFR-04 & Performance & On the documented test device/network, primary
post-login pages become usable within 3 seconds for the defined sample
dataset. \\
NFR-05 & Usability & A first-time student can complete semester,
subject, attendance, and marks workflows using the user guide without
facilitator correction in the recorded UAT. \\
NFR-06 & Accessibility & Critical forms have programmatic labels,
keyboard-reachable controls, visible focus, and no serious automated
accessibility finding in the selected audit tool. \\
NFR-07 & Compatibility & Critical journeys pass on current Chrome and
Edge desktop and one mobile viewport at 360 x 800 or a documented
equivalent. \\
NFR-08 & Maintainability & Type check, lint, automated tests, and
production build run from documented commands and pass in CI. \\
NFR-09 & Privacy & Test evidence and the repository contain no real
student records, credentials, tokens, or secret environment values. \\
NFR-10 & Recoverability & Deployment and data-change procedures state
rollback/recovery actions, ownership, and verification steps. \\
\end{longtable}
}

\section{Requirement quality
checklist}\label{144-requirement-quality-checklist}

Each approved requirement must be:

\begin{itemize}
\tightlist
\item
  Correct and necessary.
\item
  Atomic rather than combining unrelated behavior.
\item
  Unambiguous and written with \texttt{shall} for mandatory behavior.
\item
  Feasible in the five-day scope.
\item
  Verifiable through inspection, analysis, demonstration, or test.
\item
  Prioritized and uniquely identified.
\item
  Traceable backward to a need and forward to design, implementation,
  and tests.
\item
  Consistent with every other approved requirement.
\end{itemize}

\section{Critical use cases}\label{145-critical-use-cases}

At minimum, fully specify:

\begin{itemize}
\tightlist
\item
  UC-01 Register and authenticate.
\item
  UC-02 Manage a semester.
\item
  UC-03 Manage subjects and credits.
\item
  UC-04 Record attendance.
\item
  UC-05 Record examination marks.
\item
  UC-06 View SGPA/CGPA and analytics.
\item
  UC-07 Import academic data.
\item
  UC-08 Export an academic report.
\end{itemize}

Each description includes actor, trigger, preconditions, postconditions,
main flow, alternate flow, exception flow, rules, requirement IDs, and
acceptance tests.

\chapter{Analysis and Design
Standard}\label{15-analysis-and-design-standard}

\section{Architecture boundary}\label{151-architecture-boundary}

\begin{Shaded}
\begin{Highlighting}[]
\NormalTok{Student browser}
\NormalTok{  {-}\textgreater{} React pages and reusable UI components}
\NormalTok{  {-}\textgreater{} hooks, validation, domain utilities, and services}
\NormalTok{  {-}\textgreater{} Supabase client/API}
\NormalTok{  {-}\textgreater{} Supabase Authentication and PostgreSQL with RLS}
\NormalTok{  {-}\textgreater{} hosted frontend and managed backend deployment}
\end{Highlighting}
\end{Shaded}

\section{Required diagrams and
content}\label{152-required-diagrams-and-content}

{\def\LTcaptype{none} % do not increment counter
\begin{longtable}[]{@{}
  >{\raggedright\arraybackslash}p{(\linewidth - 4\tabcolsep) * \real{0.2500}}
  >{\raggedright\arraybackslash}p{(\linewidth - 4\tabcolsep) * \real{0.3846}}
  >{\raggedright\arraybackslash}p{(\linewidth - 4\tabcolsep) * \real{0.3654}}@{}}
\toprule\noalign{}
\begin{minipage}[b]{\linewidth}\raggedright
Diagram
\end{minipage} & \begin{minipage}[b]{\linewidth}\raggedright
Must show
\end{minipage} & \begin{minipage}[b]{\linewidth}\raggedright
Theory demonstrated
\end{minipage} \\
\midrule\noalign{}
\endhead
\bottomrule\noalign{}
\endlastfoot
System context & Student, UNI-TEL boundary, Supabase/auth/data, export
target, information flows & System boundary and external entities \\
Use-case & Student actor and approved user goals; relationships only
where semantically valid & Functional analysis \\
DFD Level 0 & Major processes, external entity, data stores, named data
flows & Logical process modeling \\
DFD Level 1 & Decomposition of academic record management or analytics &
Functional decomposition and balancing \\
Activity & Decisions, validation, success/failure for marks or
attendance entry & Workflow behavior \\
Sequence - login & Browser/UI, auth hook/service, Supabase Auth,
protected route & Time-ordered interaction \\
Sequence - academic record & UI, hook/service, Supabase, database,
response/error & Interface responsibility \\
Component & Pages, academic components, hooks/services, integrations,
data platform & Modular architecture \\
ERD & Profile, semester, subject, attendance, marks, keys/cardinality &
Data modeling \\
Deployment & User device, Vercel/static hosting, Supabase
Auth/API/PostgreSQL & Physical architecture \\
\end{longtable}
}

\section{Design consistency rules}\label{153-design-consistency-rules}

\begin{itemize}
\tightlist
\item
  Entity and process names must match the SRS, data dictionary, code,
  and RTM.
\item
  DFDs show data flow, not control flow or UI navigation.
\item
  Use-case diagrams show user goals, not every button.
\item
  Sequence messages correspond to real interfaces or clearly labelled
  conceptual operations.
\item
  ERD cardinality and optionality match verified foreign keys and
  constraints.
\item
  Deployment diagrams distinguish logical components from
  physical/runtime nodes.
\item
  Every significant decision has an ADR with context, alternatives,
  decision, and consequences.
\end{itemize}

\section{Database review focus}\label{154-database-review-focus}

\begin{itemize}
\tightlist
\item
  Treat migration files as authoritative until verified against a
  deployed schema.
\item
  Reconcile or remove misleading standalone schema snapshots through a
  reviewed task.
\item
  Verify primary/foreign keys, unique constraints, check constraints,
  generated values, indexes, cascades, triggers, and RLS policies.
\item
  Confirm calculations avoid division by zero and invalid numeric
  ranges.
\item
  Test with two synthetic users to demonstrate data isolation.
\item
  Document migration ordering, backup, rollback, and destructive-change
  risks.
\end{itemize}

\chapter{Implementation and Test
Strategy}\label{16-implementation-and-test-strategy}

\section{Product-hardening
principle}\label{161-product-hardening-principle}

The five-day project is not a feature expansion. Technical changes must
close a verified gap in an approved requirement, quality target,
security control, test, schema definition, accessibility review,
documentation-to-code consistency check, or release process.

\section{Planned automated test
areas}\label{162-planned-automated-test-areas}

{\def\LTcaptype{none} % do not increment counter
\begin{longtable}[]{@{}
  >{\raggedright\arraybackslash}p{(\linewidth - 4\tabcolsep) * \real{0.2157}}
  >{\raggedright\arraybackslash}p{(\linewidth - 4\tabcolsep) * \real{0.3922}}
  >{\raggedright\arraybackslash}p{(\linewidth - 4\tabcolsep) * \real{0.3922}}@{}}
\toprule\noalign{}
\begin{minipage}[b]{\linewidth}\raggedright
Area
\end{minipage} & \begin{minipage}[b]{\linewidth}\raggedright
Representative checks
\end{minipage} & \begin{minipage}[b]{\linewidth}\raggedright
Suggested location
\end{minipage} \\
\midrule\noalign{}
\endhead
\bottomrule\noalign{}
\endlastfoot
Grade calculations & Grade point mapping, weighted credits, empty
semester, invalid grade & \begin{minipage}[t]{\linewidth}\raggedright
\texttt{tests/unit/}\strut \\
\texttt{gradeCalculations.test.ts}\strut
\end{minipage} \\
Marks calculations & Valid percentage, weightage boundary, zero total
rejection, obtained greater than total &
\begin{minipage}[t]{\linewidth}\raggedright
\texttt{tests/unit/}\strut \\
\texttt{marksCalculations.test.ts}\strut
\end{minipage} \\
Attendance & Percentage, zero classes, attended greater than held,
threshold behavior & \begin{minipage}[t]{\linewidth}\raggedright
\texttt{tests/unit/}\strut \\
\texttt{attendanceCalculations.}\strut \\
\texttt{test.ts}\strut
\end{minipage} \\
Semester normalization & Normal values, imported strings, invalid
values, regression for 7 becoming 70 &
\begin{minipage}[t]{\linewidth}\raggedright
\texttt{tests/unit/}\strut \\
\texttt{semesterNormalization.}\strut \\
\texttt{test.ts}\strut
\end{minipage} \\
Validation service & Required text, numeric ranges, exam type,
cross-field validation & \begin{minipage}[t]{\linewidth}\raggedright
\texttt{tests/unit/}\strut \\
\texttt{validationService.test.ts}\strut
\end{minipage} \\
Protected routing/auth & Unauthenticated redirect, authenticated access,
loading state & \begin{minipage}[t]{\linewidth}\raggedright
\texttt{tests/integration/}\strut \\
\texttt{protectedRoute.test.tsx}\strut
\end{minipage} \\
Import & Valid fixture, malformed JSON, invalid records,
transactional/partial behavior &
\begin{minipage}[t]{\linewidth}\raggedright
\texttt{tests/integration/}\strut \\
\texttt{importAcademicData.test.ts}\strut
\end{minipage} \\
\end{longtable}
}

\section{Test levels and types}\label{163-test-levels-and-types}

\begin{itemize}
\tightlist
\item
  \textbf{Static verification:} reviews, TypeScript checks, ESLint,
  document consistency.
\item
  \textbf{Unit testing:} calculations, normalization, and validation in
  isolation.
\item
  \textbf{Integration testing:} hooks/services with Supabase boundary or
  controlled mocks; database constraints/RLS in a safe environment.
\item
  \textbf{System testing:} complete browser-level user journeys.
\item
  \textbf{Regression testing:} previously fixed semester import and all
  release-critical workflows.
\item
  \textbf{Security testing:} authentication, authorization/RLS, input
  handling, secrets, dependency review.
\item
  \textbf{Usability/accessibility testing:} task completion, navigation,
  labels, focus, contrast/audit findings.
\item
  \textbf{Compatibility testing:} supported desktop browsers and
  selected mobile viewport.
\item
  \textbf{Performance testing:} documented page/task measurements on
  defined data and environment.
\item
  \textbf{User acceptance testing:} student-perspective scenarios with
  recorded results.
\end{itemize}

\section{Test-case record}\label{164-test-case-record}

Every executed case records:

{\def\LTcaptype{none} % do not increment counter
\begin{longtable}[]{@{}
  >{\raggedright\arraybackslash}p{(\linewidth - 2\tabcolsep) * \real{0.2982}}
  >{\raggedright\arraybackslash}p{(\linewidth - 2\tabcolsep) * \real{0.7018}}@{}}
\toprule\noalign{}
\begin{minipage}[b]{\linewidth}\raggedright
Field
\end{minipage} & \begin{minipage}[b]{\linewidth}\raggedright
Example
\end{minipage} \\
\midrule\noalign{}
\endhead
\bottomrule\noalign{}
\endlastfoot
ID & TC-FR06-01 \\
Requirement & FR-06, NFR-02 \\
Objective & Verify valid marks are saved and percentage is calculated \\
Preconditions & Authenticated synthetic user; Semester 1 exists \\
Data & Subject DBMS, obtained 42, total 50, weightage 20 \\
Steps & Open Marks; add record; enter data; save; reopen record \\
Expected & Record persists; raw percentage is 84\%; weighted
contribution follows the documented formula \\
Actual & Enter the observed output after execution \\
Status & Pass, Fail, Blocked, or Not Run \\
Evidence & Screenshot/log/CI link and execution date \\
Executor/reviewer & Named member and reviewer \\
\end{longtable}
}

\section{Defect severity}\label{165-defect-severity}

{\def\LTcaptype{none} % do not increment counter
\begin{longtable}[]{@{}
  >{\raggedright\arraybackslash}p{(\linewidth - 4\tabcolsep) * \real{0.0909}}
  >{\raggedright\arraybackslash}p{(\linewidth - 4\tabcolsep) * \real{0.4545}}
  >{\raggedright\arraybackslash}p{(\linewidth - 4\tabcolsep) * \real{0.4545}}@{}}
\toprule\noalign{}
\begin{minipage}[b]{\linewidth}\raggedright
Severity
\end{minipage} & \begin{minipage}[b]{\linewidth}\raggedright
Definition
\end{minipage} & \begin{minipage}[b]{\linewidth}\raggedright
Release response
\end{minipage} \\
\midrule\noalign{}
\endhead
\bottomrule\noalign{}
\endlastfoot
Critical & Data exposure/loss, authentication bypass, build/deploy
impossible, core app unavailable & Must fix before release \\
High & Must-have workflow fails or calculation is materially wrong with
no safe workaround & Fix or formally stop release \\
Medium & Important behavior degraded with a usable workaround & Fix if
time permits; document residual risk \\
Low & Cosmetic, wording, or minor usability issue & Record for
maintenance backlog \\
\end{longtable}
}

\section{Quality commands}\label{166-quality-commands}

The DevOps and QA leads must make these concepts executable from
\texttt{package.json} and CI:

\begin{Shaded}
\begin{Highlighting}[]
\NormalTok{npm ci}
\NormalTok{npm run type{-}check}
\NormalTok{npm run lint}
\NormalTok{npm run test }\OperatorTok{{-}{-}} \OperatorTok{{-}{-}}\NormalTok{run}
\NormalTok{npm run build}
\end{Highlighting}
\end{Shaded}

Record the exact tool versions, runner OS, command output, date, and
commit SHA in the Test Report.

\chapter{Traceability Strategy}\label{17-traceability-strategy}

The RTM is the central proof that theory is displayed through the
application.

{\def\LTcaptype{none} % do not increment counter
\begin{longtable}[]{@{}
  >{\raggedright\arraybackslash}p{(\linewidth - 10\tabcolsep) * \real{0.0759}}
  >{\raggedright\arraybackslash}p{(\linewidth - 10\tabcolsep) * \real{0.1931}}
  >{\raggedright\arraybackslash}p{(\linewidth - 10\tabcolsep) * \real{0.2276}}
  >{\raggedright\arraybackslash}p{(\linewidth - 10\tabcolsep) * \real{0.2759}}
  >{\raggedright\arraybackslash}p{(\linewidth - 10\tabcolsep) * \real{0.0897}}
  >{\raggedright\arraybackslash}p{(\linewidth - 10\tabcolsep) * \real{0.1379}}@{}}
\toprule\noalign{}
\begin{minipage}[b]{\linewidth}\raggedright
Requirement
\end{minipage} & \begin{minipage}[b]{\linewidth}\raggedright
Need/use case
\end{minipage} & \begin{minipage}[b]{\linewidth}\raggedright
Design
\end{minipage} & \begin{minipage}[b]{\linewidth}\raggedright
Implementation
\end{minipage} & \begin{minipage}[b]{\linewidth}\raggedright
Test
\end{minipage} & \begin{minipage}[b]{\linewidth}\raggedright
Result/evidence
\end{minipage} \\
\midrule\noalign{}
\endhead
\bottomrule\noalign{}
\endlastfoot
FR-01 & UC-01 & Login sequence, auth architecture & \texttt{AuthPage},
\texttt{useAuth}, \texttt{ProtectedRoute} & TC-FR01-01/02 & Executed
result link \\
FR-03 & UC-02 & Semester activity/ERD & Semester page/hook/service,
\texttt{semesters} table & TC-FR03-01/02 & Executed result link \\
FR-05 & UC-04 & Attendance activity/ERD & Attendance editor/service,
\texttt{attendance\_records} & TC-FR05-01/02 & Executed result link \\
FR-06 & UC-05 & Marks sequence/ERD & Marks editor/service,
\texttt{marks\_records} & TC-FR06-01/02 & Executed result link \\
FR-08 & UC-06 & Calculation design & Grade utilities, semester
aggregation & TC-FR08-01/02 & Executed result link \\
NFR-01 & UC-01 and all data use cases & Security/RLS design & Protected
routes and RLS policies & TC-SEC-01/02 & Executed result link \\
\end{longtable}
}

Before release, the team must calculate:

\begin{itemize}
\tightlist
\item
  Requirements coverage = requirements with executed tests / approved
  requirements x 100.
\item
  Must-have pass rate = passed Must-have tests / executed Must-have
  tests x 100.
\item
  Defect closure = closed defects / recorded defects x 100.
\item
  Traceability completeness = requirements with need, design,
  implementation, test, and result links / approved requirements x 100.
\end{itemize}

Targets: 100\% traceability and executed-test coverage for Must
requirements; 100\% closure of Critical and High release-blocking
defects.

\chapter{Project Estimation and Resource
Plan}\label{18-project-estimation-and-resource-plan}

Gross capacity is 5 members x 5 days x 8 hours = 200 person-hours.
Reserve time for coordination and unexpected work.

{\def\LTcaptype{none} % do not increment counter
\begin{longtable}[]{@{}
  >{\raggedright\arraybackslash}p{(\linewidth - 4\tabcolsep) * \real{0.6452}}
  >{\raggedleft\arraybackslash}p{(\linewidth - 4\tabcolsep) * \real{0.1935}}
  >{\raggedleft\arraybackslash}p{(\linewidth - 4\tabcolsep) * \real{0.1613}}@{}}
\toprule\noalign{}
\begin{minipage}[b]{\linewidth}\raggedright
Activity
\end{minipage} & \begin{minipage}[b]{\linewidth}\raggedleft
Person-hours
\end{minipage} & \begin{minipage}[b]{\linewidth}\raggedleft
Percentage
\end{minipage} \\
\midrule\noalign{}
\endhead
\bottomrule\noalign{}
\endlastfoot
Initiation, governance, feasibility, planning & 20 & 10\% \\
Requirements and product documentation & 32 & 16\% \\
Analysis, architecture, database, diagrams & 36 & 18\% \\
Targeted implementation and automation & 40 & 20\% \\
Testing, defects, traceability, quality & 36 & 18\% \\
Deployment, release, maintenance docs & 16 & 8\% \\
Final report, slides, demo, viva & 20 & 10\% \\
\textbf{Total} & \textbf{200} & \textbf{100\%} \\
\end{longtable}
}

This is a planning allocation. Record actual effort daily without
fabricating precision. If capacity is lost, reduce Should/Could items
before compromising Must-have testing or traceability.

\chapter{Risk Management}\label{19-risk-management}

Score Probability (P) and Impact (I) from 1 to 5. Exposure = P x I.

{\def\LTcaptype{none} % do not increment counter
\begin{longtable}[]{@{}
  >{\raggedright\arraybackslash}p{(\linewidth - 12\tabcolsep) * \real{0.0609}}
  >{\raggedright\arraybackslash}p{(\linewidth - 12\tabcolsep) * \real{0.3478}}
  >{\raggedleft\arraybackslash}p{(\linewidth - 12\tabcolsep) * \real{0.0609}}
  >{\raggedleft\arraybackslash}p{(\linewidth - 12\tabcolsep) * \real{0.0609}}
  >{\raggedleft\arraybackslash}p{(\linewidth - 12\tabcolsep) * \real{0.0609}}
  >{\raggedright\arraybackslash}p{(\linewidth - 12\tabcolsep) * \real{0.3478}}
  >{\raggedright\arraybackslash}p{(\linewidth - 12\tabcolsep) * \real{0.0609}}@{}}
\toprule\noalign{}
\begin{minipage}[b]{\linewidth}\raggedright
ID
\end{minipage} & \begin{minipage}[b]{\linewidth}\raggedright
Risk
\end{minipage} & \begin{minipage}[b]{\linewidth}\raggedleft
P
\end{minipage} & \begin{minipage}[b]{\linewidth}\raggedleft
I
\end{minipage} & \begin{minipage}[b]{\linewidth}\raggedleft
Score
\end{minipage} & \begin{minipage}[b]{\linewidth}\raggedright
Response and trigger
\end{minipage} & \begin{minipage}[b]{\linewidth}\raggedright
Owner
\end{minipage} \\
\midrule\noalign{}
\endhead
\bottomrule\noalign{}
\endlastfoot
RISK-01 & Five-day schedule causes incomplete artifacts & 4 & 5 & 20 &
Freeze scope; daily gates; trigger if any Must item slips by one day &
M1 \\
RISK-02 & Artificial-looking repository history harms credibility & 3 &
5 & 15 & Transparent baseline; genuine issues/PRs; trigger on request to
backdate/split old work & M1/M5 \\
RISK-03 & Existing code and documentation disagree & 4 & 4 & 16 &
Baseline audit; traceability review; trigger on unsupported README claim
& M2/M3 \\
RISK-04 & Schema snapshots are duplicated or inconsistent & 4 & 5 & 20 &
Use migrations as authority; verify deployed schema; trigger on ERD
mismatch & M3 \\
RISK-05 & Missing automated tests allow calculation regression & 4 & 5 &
20 & Prioritize calculation/validation tests; trigger on changed domain
utility & M4 \\
RISK-06 & Supabase or hosting access fails & 3 & 4 & 12 & Local build,
screenshots, recorded backup demo, documented recovery & M5 \\
RISK-07 & Secrets or real student data enter Git/evidence & 2 & 5 & 10 &
Synthetic data, secret scan, \texttt{.gitignore}, repository audit
before release & M5/M4 \\
RISK-08 & Merge conflicts and unstable main branch & 3 & 4 & 12 & Small
PRs, WIP limit, integration checkpoints, branch protection & M5 \\
RISK-09 & Unequal contribution or review bottleneck & 3 & 4 & 12 & RACI,
review rotation, contribution dashboard, daily rebalance & M1 \\
RISK-10 & Live demo network or account failure & 3 & 5 & 15 & Seeded
demo account, local/recorded backup, screenshots, rehearsal & M2/M5 \\
RISK-11 & Documentation consumes all technical time & 3 & 4 & 12 & Reuse
verified evidence, parallel roles, YAGNI, fixed artifact outlines &
M1 \\
RISK-12 & LaTeX/report compilation fails near deadline & 2 & 3 & 6 &
Compile daily, keep PDF snapshot, freeze formatting early & M5 \\
\end{longtable}
}

Review risks at the daily retrospective. Record status, mitigation work,
changes in score, and residual risk at release.

\chapter{Change Management}\label{20-change-management}

After the Day 2 requirements baseline, every scope or requirement change
uses a Change Request containing:

\begin{itemize}
\tightlist
\item
  Change ID and requester.
\item
  Date and reason.
\item
  Affected requirement, design, code, test, document, schedule, and
  risk.
\item
  Alternatives considered.
\item
  Effort and dependency impact.
\item
  Decision: Approved, Rejected, or Deferred.
\item
  Approver and decision date.
\item
  Links to issues, PRs, tests, and RTM updates.
\end{itemize}

Approval rule:

\begin{itemize}
\tightlist
\item
  Defect correction that restores an approved requirement: Product Lead
  and relevant technical lead.
\item
  New feature or changed requirement: Product Lead plus affected leads;
  it must not endanger Must items.
\item
  Database/security/release change: mandatory QA and DevOps review.
\item
  Could-have change after Day 3: defer to maintenance backlog.
\end{itemize}

\chapter{Quality Assurance Plan}\label{21-quality-assurance-plan}

\section{Review checklists}\label{211-review-checklists}

\textbf{Requirements review:} complete, correct, feasible, singular,
unambiguous, testable, prioritized, traceable.

\textbf{Design review:} requirements coverage, clear boundaries,
consistent names, valid data relationships, security/error paths,
deployability, testability.

\textbf{Code/PR review:} linked issue, focused diff, correctness,
validation, security, tests, types, error handling, no secrets,
documentation and RTM updated.

\textbf{Test review:} requirement link, deterministic setup,
positive/negative/boundary coverage, actual result, evidence,
reproducibility, defect link for failures.

\textbf{Release review:} green CI, clean build, smoke/regression pass,
no blocking defect, complete RTM, current docs, recoverability, final
archive.

\section{Project metrics}\label{212-project-metrics}

Report values with context, not vanity numbers:

\begin{itemize}
\tightlist
\item
  Approved requirements by priority.
\item
  Requirements with design/test/result traceability.
\item
  Test cases planned/executed/passed/failed/blocked.
\item
  Automated test count and meaningful coverage areas.
\item
  Defects by severity, state, and mean resolution time where useful.
\item
  Pull requests opened/approved/merged and substantive reviews.
\item
  Build success and release-candidate status.
\item
  Planned versus completed backlog items.
\item
  Contribution evidence by category, not commits alone.
\end{itemize}

\chapter{Deployment, Operations, and
Maintenance}\label{22-deployment-operations-and-maintenance}

\section{Environment model}\label{221-environment-model}

{\def\LTcaptype{none} % do not increment counter
\begin{longtable}[]{@{}
  >{\raggedright\arraybackslash}p{(\linewidth - 6\tabcolsep) * \real{0.1453}}
  >{\raggedright\arraybackslash}p{(\linewidth - 6\tabcolsep) * \real{0.3419}}
  >{\raggedright\arraybackslash}p{(\linewidth - 6\tabcolsep) * \real{0.2479}}
  >{\raggedright\arraybackslash}p{(\linewidth - 6\tabcolsep) * \real{0.2650}}@{}}
\toprule\noalign{}
\begin{minipage}[b]{\linewidth}\raggedright
Environment
\end{minipage} & \begin{minipage}[b]{\linewidth}\raggedright
Purpose
\end{minipage} & \begin{minipage}[b]{\linewidth}\raggedright
Data
\end{minipage} & \begin{minipage}[b]{\linewidth}\raggedright
Change control
\end{minipage} \\
\midrule\noalign{}
\endhead
\bottomrule\noalign{}
\endlastfoot
Local development & Individual coding and unit testing & Synthetic/local
test data & Feature branches \\
Test/staging & Integration, RLS, UAT, release candidate & Synthetic
shared test data & Reviewed \texttt{main} or release branch \\
Production/demo & Professor demonstration and final evidence & Seeded
non-personal demo data & Approved release only \\
\end{longtable}
}

\section{Deployment guide must
include}\label{222-deployment-guide-must-include}

\begin{itemize}
\tightlist
\item
  Prerequisites and supported Node/npm versions.
\item
  Clean clone and \texttt{npm\ ci} procedure.
\item
  Environment-variable names with safe example values.
\item
  Supabase migration ordering and verification.
\item
  Local build, preview, and smoke-test commands.
\item
  Hosting configuration and deployment steps.
\item
  Release approval, version/tag, and release notes.
\item
  Rollback procedure for frontend and database changes.
\item
  Backup/recovery and data-retention assumptions.
\item
  Known limitations, service dependencies, and troubleshooting.
\end{itemize}

\section{Maintenance categories}\label{223-maintenance-categories}

\begin{itemize}
\tightlist
\item
  \textbf{Corrective:} fix defects found after release.
\item
  \textbf{Adaptive:} respond to browser, dependency, Supabase, or
  hosting changes.
\item
  \textbf{Perfective:} improve usability, performance, or supported
  reporting.
\item
  \textbf{Preventive:} dependency updates, refactoring, additional
  tests, documentation refresh.
\end{itemize}

The maintenance backlog should contain deferred Medium/Low defects and
out-of-scope enhancements without presenting them as completed release
functionality.

\chapter{Final Report Structure}\label{23-final-report-structure}

\begin{enumerate}
\def\labelenumi{\arabic{enumi}.}
\tightlist
\item
  Cover page, certificate/declaration if required, acknowledgements.
\item
  Abstract and executive summary.
\item
  Problem background, objectives, stakeholders, scope, limitations.
\item
  Feasibility study.
\item
  SDLC and methodology theory comparison.
\item
  Selected RAD/iterative/Kanban/V-Model/DevOps approach and five-day
  evidence.
\item
  Project planning: WBS, schedule, estimates, RACI, communication,
  risks.
\item
  Requirements engineering: PRD, SRS, actors, use cases, FRs/NFRs,
  priorities.
\item
  Analysis and design: architecture, DFD/UML, database, UI, security,
  decisions.
\item
  Implementation overview: modules, technologies, targeted five-day
  improvements.
\item
  Verification and validation: strategy, cases, automation, results,
  defects, RTM.
\item
  Configuration management, Git collaboration, CI/CD, release and
  deployment.
\item
  Maintenance, limitations, ethical/privacy considerations, future
  scope.
\item
  Results against success criteria and lessons learned.
\item
  Conclusion.
\item
  References using one consistent citation style.
\item
  Appendices: plans, minutes, decisions, test evidence, contribution
  matrix, user guide.
\end{enumerate}

Every theory chapter must end with a short
\texttt{Application\ to\ UNI-TEL} subsection containing real repository
evidence.

\chapter{Presentation and Demonstration
Plan}\label{24-presentation-and-demonstration-plan}

\section{Recommended 18-slide deck}\label{241-recommended-18-slide-deck}

{\def\LTcaptype{none} % do not increment counter
\begin{longtable}[]{@{}
  >{\raggedright\arraybackslash}p{(\linewidth - 4\tabcolsep) * \real{0.0972}}
  >{\raggedright\arraybackslash}p{(\linewidth - 4\tabcolsep) * \real{0.5556}}
  >{\raggedright\arraybackslash}p{(\linewidth - 4\tabcolsep) * \real{0.3472}}@{}}
\toprule\noalign{}
\begin{minipage}[b]{\linewidth}\raggedright
Slides
\end{minipage} & \begin{minipage}[b]{\linewidth}\raggedright
Subject
\end{minipage} & \begin{minipage}[b]{\linewidth}\raggedright
Speaker
\end{minipage} \\
\midrule\noalign{}
\endhead
\bottomrule\noalign{}
\endlastfoot
1-3 & Title, problem, objectives, stakeholders, scope & M1 \\
4-5 & Methodology comparison, selection, five-day SDLC & M1 \\
6-8 & PRD/SRS, FR/NFR examples, use cases, priorities & M2 \\
9-11 & Architecture, DFD/UML, ERD/database/security & M3 \\
12-14 & Test strategy, defects, automation, RTM and metrics & M4 \\
15-16 & Git collaboration, CI/CD, deployment, maintenance & M5 \\
17 & Live demonstration with speaker hand-offs & All \\
18 & Results, limitations, future scope, conclusion & M1 plus all for
questions \\
\end{longtable}
}

\section{Theory-to-evidence rule}\label{242-theory-to-evidence-rule}

Use the pattern:

\begin{Shaded}
\begin{Highlighting}[]
\NormalTok{Theory concept {-}\textgreater{} UNI{-}TEL decision/artifact {-}\textgreater{} implementation evidence {-}\textgreater{} test/result}
\end{Highlighting}
\end{Shaded}

Example:

\begin{Shaded}
\begin{Highlighting}[]
\NormalTok{Requirements traceability {-}\textgreater{} FR{-}06 {-}\textgreater{} MarksEditor + marks\_records {-}\textgreater{} TC{-}FR06{-}01 Pass}
\end{Highlighting}
\end{Shaded}

\section{Timed demo script}\label{243-timed-demo-script}

\begin{enumerate}
\def\labelenumi{\arabic{enumi}.}
\tightlist
\item
  Show repository board, requirement ID, PR, CI, and RTM row (45
  seconds).
\item
  Open the deployed landing page and authenticate (30 seconds).
\item
  Create or open a semester and show subjects/credits (45 seconds).
\item
  Enter attendance, including one validation example (45 seconds).
\item
  Enter marks and show calculated percentage/weightage (45 seconds).
\item
  Show SGPA/CGPA and analytics (45 seconds).
\item
  Export a supported report (30 seconds).
\item
  Explain RLS/user isolation and show corresponding test evidence (45
  seconds).
\item
  Show final CI/release result and state limitations (30 seconds).
\end{enumerate}

Prepare a seeded synthetic account, stable browser tabs, backup
screenshots, and a short recorded demo in case network services fail.

\section{Viva coverage by every
member}\label{244-viva-coverage-by-every-member}

Every member must be able to explain:

\begin{itemize}
\tightlist
\item
  Why RAD was selected and why Waterfall, Spiral, or pure Scrum was not.
\item
  Difference between verification and validation.
\item
  Difference between functional and non-functional requirements.
\item
  How one UNI-TEL requirement is traced to design, code, test, and
  result.
\item
  How authentication differs from authorization and how RLS helps.
\item
  Difference between unit, integration, system, regression, and
  acceptance tests.
\item
  How Git branches, PR review, CI, release, rollback, and maintenance
  work.
\item
  A limitation, a risk, a defect, a design trade-off, and a lesson
  learned.
\end{itemize}

\chapter{Evidence and Contribution
Management}\label{25-evidence-and-contribution-management}

\section{Daily evidence folders}\label{251-daily-evidence-folders}

Each day folder should contain:

\begin{itemize}
\tightlist
\item
  Stand-up and review minutes.
\item
  Board snapshot or exported issue list.
\item
  Document and diagram review evidence.
\item
  Relevant build/test output.
\item
  PR and review links.
\item
  Screenshots using synthetic data.
\item
  Updated risk, decision, defect, and contribution records.
\item
  Daily progress summary: planned, completed, carried, blocked,
  decision.
\end{itemize}

\section{Contribution matrix}\label{252-contribution-matrix}

Record quality and type, not just quantity:

{\def\LTcaptype{none} % do not increment counter
\begin{longtable}[]{@{}
  >{\raggedright\arraybackslash}p{(\linewidth - 12\tabcolsep) * \real{0.0569}}
  >{\raggedright\arraybackslash}p{(\linewidth - 12\tabcolsep) * \real{0.1138}}
  >{\raggedright\arraybackslash}p{(\linewidth - 12\tabcolsep) * \real{0.2195}}
  >{\raggedright\arraybackslash}p{(\linewidth - 12\tabcolsep) * \real{0.2033}}
  >{\raggedright\arraybackslash}p{(\linewidth - 12\tabcolsep) * \real{0.2114}}
  >{\raggedright\arraybackslash}p{(\linewidth - 12\tabcolsep) * \real{0.0813}}
  >{\raggedright\arraybackslash}p{(\linewidth - 12\tabcolsep) * \real{0.1138}}@{}}
\toprule\noalign{}
\begin{minipage}[b]{\linewidth}\raggedright
Member
\end{minipage} & \begin{minipage}[b]{\linewidth}\raggedright
Owned issues
\end{minipage} & \begin{minipage}[b]{\linewidth}\raggedright
Technical changes
\end{minipage} & \begin{minipage}[b]{\linewidth}\raggedright
Authored artifacts
\end{minipage} & \begin{minipage}[b]{\linewidth}\raggedright
Tests/evidence
\end{minipage} & \begin{minipage}[b]{\linewidth}\raggedright
PR reviews
\end{minipage} & \begin{minipage}[b]{\linewidth}\raggedright
Presentation
\end{minipage} \\
\midrule\noalign{}
\endhead
\bottomrule\noalign{}
\endlastfoot
M1 & WP1/WP7 items & Auth/security verification & Charter, PRD,
feasibility & Acceptance evidence & Reviews M5 & Slides 1-5, 18 \\
M2 & WP2/UI items & Validation/accessibility & SRS, use cases, UI guide
& UX evidence & Reviews M1 & Slides 6-8 \\
M3 & WP3/data items & Schema/integrity & SDD, diagrams, dictionary & DB
evidence & Reviews M2 & Slides 9-11 \\
M4 & WP5 items & Test framework/defect fixes & Test docs, RTM & Test and
security evidence & Reviews M3 & Slides 12-14 \\
M5 & WP0/WP6 items & CI/deployment & SCM, deploy, maintenance &
Build/release evidence & Reviews M4 & Slides 15-16 \\
\end{longtable}
}

Update with actual issue, PR, and commit links. Never assign authorship
to someone who did not perform the work.

\chapter{Final Submission
Structure}\label{26-final-submission-structure}

\begin{Shaded}
\begin{Highlighting}[]
\NormalTok{submission/}
\NormalTok{|{-}{-} 01\_Project\_Charter\_and\_Proposal.pdf}
\NormalTok{|{-}{-} 02\_Feasibility\_and\_Methodology.pdf}
\NormalTok{|{-}{-} 03\_Project\_Management\_Plan.pdf}
\NormalTok{|{-}{-} 04\_PRD.pdf}
\NormalTok{|{-}{-} 05\_SRS\_and\_Use\_Cases.pdf}
\NormalTok{|{-}{-} 06\_Software\_Design\_and\_Diagrams.pdf}
\NormalTok{|{-}{-} 07\_Database\_Design.pdf}
\NormalTok{|{-}{-} 08\_Test\_Plan\_Cases\_and\_Report.pdf}
\NormalTok{|{-}{-} 09\_Requirements\_Traceability\_Matrix.xlsx{-}or{-}pdf}
\NormalTok{|{-}{-} 10\_Risk\_Change\_and\_Configuration\_Management.pdf}
\NormalTok{|{-}{-} 11\_Deployment\_Maintenance\_and\_User\_Guide.pdf}
\NormalTok{|{-}{-} 12\_Final\_Project\_Report.pdf}
\NormalTok{|{-}{-} 13\_Presentation.pdf}
\NormalTok{|{-}{-} 14\_Team\_Contribution\_and\_Evidence.pdf}
\NormalTok{|{-}{-} source{-}code{-}and{-}document{-}sources/}
\NormalTok{|{-}{-} release{-}notes/}
\NormalTok{\textasciigrave{}{-}{-} README\_Submission\_Index.pdf}
\end{Highlighting}
\end{Shaded}

Do not duplicate contradictory content across documents. The Submission
Index identifies the authoritative version and location of every
artifact.

\chapter{Final Audit Checklists}\label{27-final-audit-checklists}

\section{SE completeness}\label{271-se-completeness}

\begin{itemize}
\tightlist
\item[$\square$]
  Problem, vision, objectives, stakeholders, scope, exclusions,
  assumptions, constraints.
\item[$\square$]
  Technical, operational, economic, schedule, legal/security
  feasibility.
\item[$\square$]
  Methodology theory comparison and justified selection.
\item[$\square$]
  WBS, schedule/Gantt, estimates, roles, RACI, communications,
  governance.
\item[$\square$]
  PRD, SRS, actors, use cases, FRs, NFRs, rules, priorities, acceptance
  criteria.
\item[$\square$]
  Context, use-case, DFD 0/1, activity, sequence, component, ER,
  deployment diagrams.
\item[$\square$]
  Architecture, modules, interfaces, database, UI, validation, errors,
  security design.
\item[$\square$]
  Implementation evidence tied to approved issues and requirements.
\item[$\square$]
  Static, unit, integration, system, regression, security, UX,
  compatibility, UAT evidence.
\item[$\square$]
  Test Plan, cases/results, Defect Log, Test Report, and complete RTM.
\item[$\square$]
  Risk, change, configuration, version, review, and quality management.
\item[$\square$]
  CI/CD, environment, deployment, rollback, recovery, maintenance, and
  user guide.
\item[$\square$]
  Final report, presentation, demo, viva, contribution record, and
  release archive.
\end{itemize}

\section{Repository and release}\label{272-repository-and-release}

\begin{itemize}
\tightlist
\item[$\square$]
  Imported baseline is transparent and old \texttt{.git} history is
  absent.
\item[$\square$]
  No credentials, \texttt{.env}, personal data, build output, or
  temporary evidence is tracked.
\item[$\square$]
  Issues, branches, commits, PRs, and reviews are genuine and linked.
\item[$\square$]
  \texttt{npm\ ci}, type check, lint, automated tests, and build pass on
  the final SHA.
\item[$\square$]
  Database migration and RLS evidence are verified safely.
\item[$\square$]
  All Critical/High release blockers are resolved or release is stopped.
\item[$\square$]
  RTM covers all Must requirements with executed results.
\item[$\square$]
  Clean deployment and smoke test are successful.
\item[$\square$]
  Tag \texttt{v1.0.0-se-class} and release notes exist.
\item[$\square$]
  Submission archive can be opened and understood on another machine.
\end{itemize}

\section{Presentation readiness}\label{273-presentation-readiness}

\begin{itemize}
\tightlist
\item[$\square$]
  Every theory claim has project evidence.
\item[$\square$]
  All five members have balanced speaking roles and can answer common
  viva questions.
\item[$\square$]
  Diagrams and screenshots are legible from presentation distance.
\item[$\square$]
  Demo account uses synthetic data and has been reset.
\item[$\square$]
  Live demo is rehearsed and timed.
\item[$\square$]
  Offline screenshots/video and local build are available as backup.
\item[$\square$]
  Limitations and future scope are stated honestly.
\end{itemize}

\chapter{Reusable Record Templates}\label{28-reusable-record-templates}

\section{Daily progress record}\label{281-daily-progress-record}

\begin{Shaded}
\begin{Highlighting}[]
\NormalTok{Date: 14 August 2026}
\NormalTok{Iteration goal: Establish verified baseline, scope, governance, and approved Day 2 work.}
\NormalTok{Completed: List issue/PR/document links.}
\NormalTok{Carried forward: State item, reason, owner, and new due date.}
\NormalTok{Blocked: State blocker, impact, owner, and escalation.}
\NormalTok{Decisions: Link ADR or meeting decision.}
\NormalTok{Risks changed: Link Risk Register entries.}
\NormalTok{Quality result: Record build/test/document gate.}
\NormalTok{Approvals: Project Lead and rotating reviewer.}
\end{Highlighting}
\end{Shaded}

\section{Meeting minutes record}\label{282-meeting-minutes-record}

\begin{Shaded}
\begin{Highlighting}[]
\NormalTok{Meeting: Day 2 Requirements Review}
\NormalTok{Date/time: 15 August 2026, 17:30 IST}
\NormalTok{Attendees: Members 1{-}5}
\NormalTok{Inputs: PRD v1.0, SRS review candidate, Use Cases, Test Plan v0.1}
\NormalTok{Decisions: Record approved/rejected requirement and reason.}
\NormalTok{Actions: Record owner, due date, and issue link.}
\NormalTok{Risks/changes: Record Risk or Change Request ID.}
\NormalTok{Approval: Member 1 with reviewer confirmation.}
\end{Highlighting}
\end{Shaded}

\section{Architecture decision
record}\label{283-architecture-decision-record}

\begin{Shaded}
\begin{Highlighting}[]
\NormalTok{ADR{-}001: Use RAD as the primary five{-}day lifecycle}
\NormalTok{Status: Accepted}
\NormalTok{Context: Five calendar days, existing prototype, five{-}person team, full SE evidence required.}
\NormalTok{Options: Waterfall, Scrum, Spiral, RAD with iterative/Kanban/DevOps practices.}
\NormalTok{Decision: RAD with daily increments, Kanban flow, V{-}Model traceability, and CI.}
\NormalTok{Consequences: Fast execution and reuse; strict gates are required to prevent documentation and quality gaps.}
\end{Highlighting}
\end{Shaded}

\section{Change request example}\label{284-change-request-example}

\begin{Shaded}
\begin{Highlighting}[]
\NormalTok{CR{-}01: Add teacher dashboard}
\NormalTok{Reason: Suggested extension during Day 3.}
\NormalTok{Impact: New actor, roles, RLS policies, screens, requirements, tests, and schedule risk.}
\NormalTok{Decision: Deferred.}
\NormalTok{Justification: Outside the frozen student{-}only scope and endangers Must{-}have verification.}
\NormalTok{Target: Maintenance/future{-}scope backlog.}
\NormalTok{Approver: Project Lead after architecture and QA consultation.}
\end{Highlighting}
\end{Shaded}

\section{Defect example}\label{285-defect-example}

\begin{Shaded}
\begin{Highlighting}[]
\NormalTok{DEF{-}01: Invalid obtained marks can exceed total marks}
\NormalTok{Requirement: FR{-}06, NFR{-}02}
\NormalTok{Severity: High}
\NormalTok{Environment/commit: Record exact test environment and SHA.}
\NormalTok{Reproduction: Create marks record with obtained 60 and total 50.}
\NormalTok{Expected: Submission is rejected with a clear validation message.}
\NormalTok{Actual: Record observed behavior after execution.}
\NormalTok{Owner/reviewer: Assigned developer and QA reviewer.}
\NormalTok{Resolution/retest: Link fix PR and executed regression result.}
\end{Highlighting}
\end{Shaded}

\chapter{Immediate Start Checklist}\label{29-immediate-start-checklist}

Complete these actions in order today:

\begin{enumerate}
\def\labelenumi{\arabic{enumi}.}
\tightlist
\item
  Create the new private/public repository according to class rules.
\item
  Import the current project as the single transparent baseline commit.
\item
  Verify \texttt{.gitignore}, \texttt{.env.example}, licenses, and
  absence of secrets.
\item
  Confirm a clean \texttt{npm\ ci} and \texttt{npm\ run\ build} from the
  new repository.
\item
  Add all five collaborators and protect \texttt{main}.
\item
  Create the milestone, Kanban columns, labels, issue and PR templates.
\item
  Replace Member 1-5 labels with names and confirm review rotation.
\item
  Create Day 1 issues SE-001 through SE-006 with acceptance criteria.
\item
  Run the baseline demonstration and parallel audits.
\item
  Finish the Day 1 gate before starting requirements work.
\end{enumerate}

\chapter{Approval}\label{30-approval}

The team approves this plan when all five members agree to the scope,
methodology, roles, contribution rules, five-day schedule, evidence
policy, and release gates. Any later change must follow the Change
Management section.

{\def\LTcaptype{none} % do not increment counter
\begin{longtable}[]{@{}
  >{\raggedright\arraybackslash}p{(\linewidth - 2\tabcolsep) * \real{0.4521}}
  >{\raggedright\arraybackslash}p{(\linewidth - 2\tabcolsep) * \real{0.5479}}@{}}
\toprule\noalign{}
\begin{minipage}[b]{\linewidth}\raggedright
Role
\end{minipage} & \begin{minipage}[b]{\linewidth}\raggedright
Approval responsibility
\end{minipage} \\
\midrule\noalign{}
\endhead
\bottomrule\noalign{}
\endlastfoot
Member 1 / Project Lead & Scope, schedule, release decision \\
Member 2 / BA-UX Lead & Requirements and usability completeness \\
Member 3 / Architecture-Data Lead & Design and data correctness \\
Member 4 / QA-Security Lead & Verification, defects, and traceability \\
Member 5 / DevOps-Release Lead & Configuration, build, deployment,
archive \\
\end{longtable}
}

\begin{center}\rule{0.5\linewidth}{0.5pt}\end{center}

\textbf{Planning principle:} A credible five-day project is demonstrated
by a transparent baseline, disciplined decisions, traceable artifacts,
tested corrections, meaningful peer review, and a reproducible release.
The goal is not to imitate months of coding history; it is to show that
the team can apply the complete Software Engineering process to a real
system under a strict time box.

\end{document}
